\documentclass[11pt]{article}
\usepackage[T1]{fontenc}
\usepackage[margin=1in]{geometry}
\usepackage{amsmath,amssymb,amsthm}
\usepackage[hidelinks]{hyperref}
\newtheorem{proposition}{Proposition}[section]
\newtheorem{theorem}[proposition]{Theorem}
\newtheorem{lemma}[proposition]{Lemma}
\newtheorem{corollary}[proposition]{Corollary}
\newcommand{\T}{\mathbb T}
\newcommand{\C}{\mathbb C}
\newcommand{\dm}{\,dm}
\newcommand{\norm}[1]{\left\lVert#1\right\rVert}
\title{A counterexample to canonical Toeplitz eigenvalue distribution}
\author{Clemens Thalhammer}
\date{9 October 2026}
\begin{document}
\maketitle
\begin{abstract}
We construct a computably specified continuous complex scalar symbol
with neither an inner nor an outer annular holomorphic extension,
whose Toeplitz sections are not canonically distributed. Along an
infinite sequence of orders, each of two points outside the symbol
range is an eigenvalue of fixed positive asymptotic multiplicity.
The construction uses exact finite spectral corrections, geometric
Cauchy coordinates, and a uniformly convergent Fourier series.
A deterministic search using finite contraction certificates specifies
every correction. All parameter choices are explicit, with no claim
of computational efficiency.
\end{abstract}

\section{Statement and notation}
For a continuous complex scalar symbol $a$ on $\T$, write
\[
 a_k=\frac1{2\pi}\int_0^{2\pi}a(e^{it})e^{-ikt}\,dt,\qquad
 T_n(a)=(a_{j-k})_{j,k=0}^{n-1}.
\]
The canonical-distribution conjecture, in the continuous-symbol
formulation recorded in \cite[Section~1, p.~2]{BBG}, asserts convergence
of the empirical eigenvalue measures to $a_*m$ when both separate
one-sided annular extensions are absent. The problem is catalogued as
OP-00159 at clopen.solutions~\cite{Clopen} and as SP-14 in
OpenProblemsInNLA~\cite{OpenProblemsInNLA}.
Theorem~\ref{thm:counterexample}
gives a counterexample under exactly these hypotheses. Its symbol has the form
\[
 a(z)=zg(z^2),\qquad g_0(s)=\sqrt{1+s^{-1}},\qquad
 g=g_0+P_-+P_+.
\]
Every finite background $P_-$ has strictly negative powers and
$P_+$ has strictly positive powers; each separately vanishes at $-1$.
The square root is the branch analytic outside the unit circle,
equal to one at infinity and continuous on the circle.
Set $a_0(z)=zg_0(z^2)$ and $h(s)=sg(s)^2-1$.
Smallness of $h-s$ in the norms below makes $h(\T)$ a Jordan curve
with bounded interior $D$ containing zero.
Its normalized conformal map is $f:\mathbb D\to D$, $f(0)=0$,
$f'(0)>0$; real symmetry gives $f(-1)=-1$.

For Laurent series and analytic series, respectively, use
\[
 \|v\|_{\mathcal W^\beta}=\sum_{j\in\mathbb Z}
                 (1+|j|)^\beta|v_j|,\qquad
 \|x\|_{H^s}^2=\sum_{k\ge0}(k+1)^{2s}|x_k|^2.
\]
The two-sided version of $H^s$ uses $1+|j|$.
Write $P_nx=\sum_{k=0}^{n-1}x_kz^k$.
Write $\Pi_+,\Pi_-$ for the nonnegative and strictly negative
Fourier projections, and $T(g_0)V=\Pi_+(g_0V)$ on analytic
polynomials. The background polynomials remain denoted by $P_-,P_+$.

\section{Exact finite algebra and the boundary inverse estimates}

\begin{proposition}[Odd-symbol pencil and its jets]
\label{prop:odd-pencil}
Let $a(z)=zg(z^2)$, $n=2m+1$, $L=m+1$, and let $U$ be the
$L$-by-$L$ upper unit shift. Put
\[
 G=T_L(g),\qquad H=G^2-U,\qquad
 M(t)=e_0^*(H-tU)^{-1}e_m.
\]
Whenever $H$ is invertible, divisibility of
$\det(wI-T_n(a))/w$ by $(w^2-1)^h$ is equivalent to
$[t^k]M(t)=0$ for $0\le k<h$.
This is an exact equivalence of finite algebraic constraints.
If $W=H^{-1}$ and $r_k=e_0^*W(UW)^k$, then
\[
 [t^k]M(t)=(r_k)_m.
\]
Write $R_k(s)=\sum_{j=0}^m(r_k)_js^j$ and
$X_k(s)=\sum_{j=0}^m(r_kG)_js^j$. For a change $v_q$ in $g_{-q-1}$,
\begin{equation}\label{eq:odd-pencil-jacobian}
 \frac{\partial [t^k]M(t)}{\partial v_q}
 =-2\sum_{l=0}^k[s^{m-q-1}]R_l(s)X_{k-l}(s).
\end{equation}
At $g(s)=\sqrt{1+s^{-1}}$ one has $H=I$, $M(t)=t^m$, and
\[
 J^0_{kq}=-2(k+1)\binom{1/2}{m-q-1-k}
             \mathbf1_{k\le m-q-1}.
\]
Finally, if $\|g-\sqrt{1+s^{-1}}\|_\infty\le\eta$, then
\[
 \|H-I\|\le 2\sqrt2\eta+\eta^2.
\]
Thus sufficiently small symbol perturbations give invertibility of $H$
uniformly in $m$, without an asymptotic finite-section argument.
\end{proposition}
\begin{proof}
Let $c_w(z)=z^{-1}(a(z)-w)$ and order the even coordinates before
the odd ones in $T_{2m+2}(c_w)$. Its block form is
\[
 \begin{pmatrix}G&-wI\\-wU&G\end{pmatrix}.
\]
Its determinant is $\det(G^2-w^2U)$: first use the Schur complement
when $G$ is invertible, then use the polynomial identity to remove
that restriction. For $w\ne0$, solving for the first inverse row gives
the two blocks
\[
 e_0^*(G^2-w^2U)^{-1}G,
 \qquad we_0^*(G^2-w^2U)^{-1}.
\]
Jacobi's identity at the upper-right corner gives
\[
 \det(wI-T_n(a))
 =w\det(H-tU)M(t),\qquad t=w^2-1.
\]
The two signs cancel because $n$ is odd. Since $\det H\ne0$,
the factor multiplying $M$ is nonzero at $t=0$. Oddness of the
characteristic polynomial supplies the matching multiplicity at $w=-1$.
The resolvent series proves the row-jet formula.

All Toeplitz matrices here are persymmetric with respect to the reversal
matrix $J_L$. So are $H-tU$ and its inverse, and each coefficient
$W(UW)^k$. Their last columns are the reversals of their first rows,
with a transpose and no complex conjugation. Changing $g_{-q-1}$ changes
$G$ by $E=U^{q+1}$ and changes $H$ to first order by $EG+GE$.
Differentiating the resolvent gives
\[
 -\sum_{l=0}^k r_l(EG+GE)J_Lr_{k-l}^{\mathsf T}.
\]
The two sums agree after reversing $l$ and $k-l$. The remaining matrix
product selects the coefficient of degree $m-q-1$, proving
\eqref{eq:odd-pencil-jacobian}. At the base symbol,
$G=\sqrt{I+U}$, so $G^2-U=I$; the displayed triangular formula follows.
The multiplier bound $\|\sqrt{I+U}\|\le\sqrt2$ and
$\|G-G_0\|\le\eta$ prove the last assertion.
\end{proof}


\begin{lemma}[Changing the spectral coordinate exactly]
\label{lem:coordinate-jets}
Let $f(u)=\sum_{l\ge1}f_lu^l$ be analytic near zero with $f_1\ne0$,
and hold this function fixed while varying $G$.
The first $h$ jets of
$e_0^*(H-f(u)U)^{-1}e_m$ vanish if and only if the first $h$ jets
of $M(t)$ vanish. Set $W=H^{-1}$ and define rows
\[
 r_0=e_0^*W,\qquad
 r_k=\left(\sum_{l=1}^k f_l r_{k-l}\right)UW.
\]
Then $r_k$ is the $k$th row coefficient of the transformed inverse.
The Jacobian formula \eqref{eq:odd-pencil-jacobian} remains valid with
these rows and their products with $G$.
\end{lemma}
\begin{proof}
Composition by $f$ acts on the first $h$ Taylor coefficients through
a lower triangular matrix with diagonal $1,f_1,\ldots,f_1^{h-1}$.
This matrix is invertible. The resolvent identity gives the row
recursion, and differentiation and persymmetry give the unchanged
bilinear formula. Fixing $f$ avoids extra derivatives from the coordinate
choice; an adaptive coordinate would require accounting for them.
\end{proof}


\begin{lemma}[Suppression of repeated corner corrections]
\label{lem:corner-feedback}
Suppose $G=T_{m+1}(g)$ has lower bandwidth $p$, and set $H=G^2-U$.
Let $h+p<(m+2)/2$, $Q=m-h$, and
$E=\sum_{q=Q}^{m-1}v_qU^{q+1}$. Then $E^2=0$ and
\[
 (G+E)^2-U=H+D,\qquad D=GE+EG.
\]
Let $A$ select coordinates $0,\ldots,h+p-1$, and let $B$ select
$m+1-h-p,\ldots,m$. Then $D=ADB$.
Suppose, for $|t|\le R$, that
\[
 \|(H-tU)^{-1}\|\le M_0,\qquad
 \|D_r(H-tU)^{-1}D_r^{-1}\|\le M_r,
 \quad D_r=\operatorname{diag}(1,r,\ldots,r^m),\quad r>1.
\]
Put $d=m+2-2h-2p$ and $\gamma=M_rr^{-d}$. If $\gamma\|D\|<1$,
the nonlinear remainder in the corner after its first variation obeys
\[
 \left|[t^k]e_0^*\left((H+D-tU)^{-1}-(H-tU)^{-1}
       +(H-tU)^{-1}D(H-tU)^{-1}\right)e_m\right|
 \le R^{-k}\frac{M_0^2\gamma\|D\|^2}{1-\gamma\|D\|}.
\]
\end{lemma}
\begin{proof}
The support statements follow by multiplying the indicated shifts by a
matrix of lower bandwidth $p$. If $W=(H-tU)^{-1}$, the weighted inverse
bound gives $\|BWA\|\le M_rr^{-d}$. Each term beyond the first variation
in the resolvent identity contains this block between two copies of $D$.
Summing its geometric series gives the stated uniform bound on the disk;
Cauchy's coefficient estimate supplies $R^{-k}$. The same series proves
invertibility of the perturbed pencil on that disk.
\end{proof}


\begin{proposition}[Geometric factor and exponential separation]
\label{prop:outside-factor}
Let $c$ extend holomorphically to $|z|>1$, continuously to $\T$, and
meromorphically to infinity with pole order $b\ge0$ and nonzero leading
coefficient. Suppose $c$ is nonzero on $\T$ with winding zero.
There are exactly $b$ zeros $\zeta_1,\ldots,\zeta_b$ in $|z|>1$,
counting multiplicity. Put
\[
 p(z)=\prod_{j=1}^b(1-z/\zeta_j),\qquad
 F(z)=\frac{p(z^{-1})}{c(z^{-1})}.
\]
The apparent singularity of $F$ at zero is removable, and $F$ is the
generating function of the first row of $T(c)^{-1}$.

Fix $1<r<\min_j|\zeta_j|$ (any $r>1$ in the zero-free exterior if
$b=0$), and write $F(z)=\sum_{j\ge0}f_jz^j$.
There are constants $C,N_0$ such that, for $N\ge N_0$ and
$1\le l\le N-b$,
\begin{equation}\label{eq:low-row-exponential}
 \left(\sum_{j=0}^{l-1}|(T_N(c)^{-1})_{0j}-f_j|^2\right)^{1/2}
 \le C r^{-(N-b-l+1)}.
\end{equation}
The constants are uniform over a compact holomorphic family of such
symbols when the exterior pole order is fixed, the leading coefficient
stays nonzero, and the zeros stay outside the chosen radius $r$.
For a variation $\delta c(z)=z^{-j}$, $j\ge1$, the coefficients of
the outside-zero polynomial satisfy $\|\delta p\|=O(r^{-j})$, with
any fixed coefficient norm and the same compact-family qualification.
\end{proposition}
\begin{proof}
Nonvanishing and continuity on the boundary give a zero-free collar.
The argument principle on a large annulus gives the zero count $b$.
Choose the nonzero constant $G$ so that
$c_+=Gp$, $c_-=c/(Gp)$ satisfy $c_-(\infty)=1$. Both factors and
their inverses are bounded and holomorphic on their respective sides
of the circle. The elementary Toeplitz product identity gives
$T(c)=T(c_-)T(c_+)$, and hence
$T(c)^{-1}=T(c_+^{-1})T(c_-^{-1})$. Its first row has generating
function $G^{-1}c_-^{-1}(z^{-1})=F(z)$.

The symbols $c(rz)$ are again continuous, nonvanishing and of index
zero. Finite-section stability, in the form read in
\cite[Section~3, equation~(6)]{BFGM}, gives a uniform bound for
\[
 D_rT_N(c)^{-1}D_r^{-1}=T_N(c(r\,\cdot))^{-1},
 \qquad D_r=\operatorname{diag}(1,r,\ldots,r^{N-1}).
\]
Truncating the infinite inverse row $f$ leaves a residual supported
only in the last $b$ columns: $c_k=0$ for $k>b$, so a missing index
$i\ge N$ can contribute to column $j$ only when $j\ge N-b$.
Its norm is bounded by a constant times $\|f\|_2$.
Multiplication by the finite inverse and its displayed weighted bound
then proves \eqref{eq:low-row-exponential}. Uniformity over compact
families follows by a finite covering and the inverse perturbation
identity, applied also to the scaled symbols.

For completeness, the coefficients of $p$ depend holomorphically on a
holomorphic parameter even when the outside zeros coincide. Their
power sums are contour integrals of $z^k c'(z)/c(z)$ on an annulus
whose inner boundary is $|z|=r$ and whose outer boundary encloses all
zeros. Differentiating and integrating by parts changes the integrand
to $-kz^{k-1}\delta c(z)/c(z)$. It is $O(r^{-j})$ for the indicated
variation on the two fixed contours. Newton's identities and the
nonzero constant term of the monic zero polynomial give the final claim.
\end{proof}


\begin{lemma}[The finite corner and a squared outside factor]
\label{lem:corner-scattering-factor}
Under the hypotheses of Proposition~\ref{prop:outside-factor}, write
$p(z)=\sum_{j=0}^b p_jz^j$, $F(z)=\sum_{j\ge0}f_jz^j$,
and $V_N=T_N(c)^{-1}$. For every fixed allowable $r>1$,
\begin{equation}\label{eq:corner-scattering-factor}
 (V_N)_{0,N-1}
 =[z^{N-1}]\frac{p(z^{-1})^2}{c(z^{-1})}+O(r^{-N}).
\end{equation}
The coefficient on the right is a Laurent coefficient; its negative
part is finite. Constants are uniform on the compact families in that
proposition. No relative error is asserted.
\end{lemma}
\begin{proof}
For $N>b$, truncating the infinite first inverse row gives the exact
identity
\[
 (V_N)_{0,N-1}=f_{N-1}
 +\sum_{l=0}^{b-1}f_{N+l}
       \sum_{k=0}^{b-1-l}c_{k+l+1}(V_N)_{0k}.
\]
Indeed the omitted tail leaves a residual only in the last $b$ columns,
and persymmetry identifies $(V_N)_{N-1-k,N-1}$ with $(V_N)_{0k}$.
Equation~\eqref{eq:low-row-exponential}, with $l=b$, replaces the
finitely many entries $(V_N)_{0k}$ by $f_k$, at error $O(r^{-N})$;
the fixed factors $r^{2b-1}$ are absorbed into the constant.
The identity $F(z)c(z^{-1})=p(z^{-1})$ gives
\[
 \sum_{k=0}^{b-1-l} f_kc_{k+l+1}=p_{l+1}.
\]
The resulting sum $\sum_{j=0}^b p_jf_{N-1+j}$ is exactly the
coefficient in \eqref{eq:corner-scattering-factor}. The case $b=0$
is the exact upper-triangular identity.
\end{proof}


\begin{lemma}[Continuation of the far-zero factors across the boundary]
\label{lem:far-factor-collar}
For $g=g_0+P$ with $P(-1)=0$ and $h-s$ sufficiently small
in $C^1(\T)$, the factors $p_w$ can be chosen
holomorphically for all $w$ in a neighborhood of
$\{w:w^2-1\in\overline D\}$. There is a fixed radius $\rho>1$
such that they consist precisely of the zeros of $c_w$ outside
$|z|=\rho$ there. If the positive degree of $g$ is $p>0$, each factor
has degree $2p$; when $g$ has no positive powers the factor is one.
There are polynomials $A_t$ and $B_t$ in $s$, with coefficients
holomorphic in a neighborhood of $\overline D$, such that
\begin{equation}\label{eq:outside-square-parity}
 p_w(z)^2=A_t(z^2)+wzB_t(z^2),\qquad t=w^2-1.
\end{equation}
\end{lemma}
\begin{proof}
The function $h$ is holomorphic outside the circle, and $h'$ extends
continuously to it: the additional term $2sg_0P$ has derivative
$O(|s+1|^{1/2})$ at $-1$. For a sufficiently thin closed exterior
annulus, $\|h'-1\|<1/3$. Any two points of this annulus can be joined
in it by a radial segment and a shorter circular arc of total length
at most $(1+\pi/2)$ times their distance. Integration of $h'-1$
along that path proves injectivity. Thus $h$ maps the annulus
conformally to the region between $\Gamma$ and an enclosing Jordan
curve $h(r\T)$, for some $r>1$.

Let $U$ be the bounded interior of this latter curve. For $t\in U$,
the zeros of $h(s)-t$ outside $|s|=r$ have constant count $2p$ when
$p>0$, and zero count when $p=0$, by the argument principle at
infinity. The domain $\{w:w^2-1\in U\}$ is connected, since $U$
is simply connected and contains $-1$. On $|z|=\sqrt r$, neither
$c_w$ nor $c_{-w}$ vanishes there. Their equal outside-zero counts
are $2p$: at $w=0$ this follows from
$g(s)^2=(h(s)+1)/s$ and the argument principle. Choosing
$\rho=\sqrt r$, contour power sums therefore give holomorphic
normalized factors, including at multiple far zeros. For $t\in D$
they agree with the original factors, because the intervening
annulus has no zeros. The parity relation $p_{-w}(z)=p_w(-z)$ shows
that the even coefficients in $p_w^2$ are even in $w$, while its
odd coefficients are odd in $w$. Division of the latter by $w$ is
removable at zero. Descent under $w\mapsto w^2-1$ proves
\eqref{eq:outside-square-parity}.
\end{proof}



For the factors in the preceding lemma define $R_t,X_t$ by parity:
\[
 F_w(z)=\frac{p_w(z^{-1})}{g(z^{-2})-wz}
          =X_t(z^2)+wzR_t(z^2),\qquad t=w^2-1.
\]
Parity and the removable value at $w=0$ make both functions
holomorphic in $t$. Their coefficients are the formal inverse rows
used below.

\begin{proposition}[Polynomial stability up to the squared boundary]
\label{prop:boundary-polynomial-stability}
For the fixed background of Lemma~\ref{lem:far-factor-collar}, put
$\mathcal W=\{w:w^2-1\in\overline D\}$.
There are $C,N_0$ and $\rho>1$ such that, for every $w\in\mathcal W$
and $N\ge N_0$, $T_N(c_w)$ is invertible and
\[
 \|T_N(c_w)^{-1}\|\le CN^2,\qquad
 \|D_\rho T_N(c_w)^{-1}D_\rho^{-1}\|\le C,
 \quad D_\rho=\operatorname{diag}(1,\rho,\ldots,\rho^{N-1}).
\]
These include the parameters for which $c_w$ vanishes on $\T$.
For the odd pencil $W_m(t)=(G^2-(1+t)U)^{-1}$ it follows that,
uniformly for $t\in\overline D$ and large $m$,
\begin{equation}\label{eq:boundary-pencil-stability}
 \|W_m(t)\|\le C(m+1)^4,\qquad
 \|D_rW_m(t)D_r^{-1}\|\le C,\qquad r=\rho^2>1.
\end{equation}
\end{proposition}
\begin{proof}
The scaled symbols $c_w(\rho\,\cdot)$ are nonzero and of index zero
uniformly on the compact parameter set $\mathcal W$, by the collar
lemma. Uniform finite-section stability gives the second bound in
the first display, and also eventual invertibility. This step applies
to the scaled symbols, not to a possibly vanishing symbol on $\T$.
If $V_N=T_N(c_w)^{-1}$, in particular
\[
 |(V_N)_{ij}|\le C\rho^{j-i}.
\]

The far-factor function $F_w(z)=p_w(z^{-1})/c_w(z^{-1})$ is
holomorphic for $|z|<1$, also when $w\in\partial\mathcal W$.
Cancellation removes the far zeros. Near the circle, rationalization
of its denominator expresses it as
\[
 F_w(z)=\frac{p_w(z^{-1})(g(z^{-2})+wz)}
                {z^2(h(z^{-2})-(w^2-1))}.
\]
The collar map is quantitatively injective. Since $h(s)$ is outside
$D$ when $1<|s|<r$, this gives
$\operatorname{dist}(h(s),\overline D)\ge c(|s|-1)$.
The numerator is uniformly bounded there. Hence
\[
 |F_w(z)|\le\frac C{1-|z|},\qquad
 |[z^j]F_w(z)|\le C(j+1),
\]
uniformly in $w$; the second inequality is Cauchy's bound at
radius $j/(j+1)$, with the small values of $j$ treated separately.

The formal identity $F_w(z)c_w(z^{-1})=p_w(z^{-1})$ still holds at
boundary parameters. Its row convolution is finite in each column
because $c_w$ has fixed positive degree $b$. Truncation therefore
gives a residual of size at most $CN$ supported in the last $b$
columns. Multiplication by the displayed entry bound for $V_N$ shows
\[
 |(V_N)_{0j}-[z^j]F_w|\le CN\rho^{j-N+b}.
\]
In particular, the first row has entries bounded by $CN$ and
$(V_N)_{00}=F_w(0)+O(N\rho^{-N})$. The continuous function $F_w(0)$
never vanishes: it is the ratio of the leading coefficients of the
normalized far factor and $c_w$. Thus $|(V_N)_{00}|$ is bounded
below for large $N$.

For clarity, the elementary Toeplitz inverse identity needed to pass
from a row bound to an operator bound is as follows. Write
$x=V_Ne_0$, $y=V_Ne_{N-1}$, and $x_0=(V_N)_{00}$.
If $L(v)$ and $U(v)$ denote lower and upper triangular Toeplitz
matrices with first column or row $v$, respectively, then
\[
 V_N=\frac1{x_0}\left[
 L(x)U(y_{N-1},\ldots,y_0)
 -L(0,y_0,\ldots,y_{N-2})U(0,x_{N-1},\ldots,x_1)\right].
\]
One verification uses the equal first and last principal blocks of
the original Toeplitz matrix. Their inverse Schur complements give
\[
 x_0\{(V_N)_{ij}-(V_N)_{i-1,j-1}\}
 =x_i y_{N-1-j}-y_{i-1}x_{N-j},\qquad i,j\ge1;
\]
iteration with the first row and column gives the identity. The
Schur complements exist since $x_0\ne0$ and the matrix is
persymmetric. Now $\|x\|_1\le C$ by the weighted entry estimate,
whereas $\|y\|_1\le CN^2$ by persymmetry and the first-row bound.
The triangle and convolution inequalities prove $\|V_N\|\le CN^2$.

Finally use the even--odd block matrix in
Proposition~\ref{prop:odd-pencil}. At $w=0$ its inverse bounds give
$\|G^{-1}\|\le C(m+1)^2$. Its upper-left inverse block is $W_m(t)G$,
so multiplication by $G^{-1}$ proves the unweighted pencil bound.
After scaling by $\rho$ in the original variable, the blocks use
$G_r=T_{m+1}(g(r\,\cdot))$, and the off-diagonal factors are
$w/\rho$. Both that full inverse and $G_r^{-1}$ have uniform bounds.
The same identity proves the weighted bound in
\eqref{eq:boundary-pencil-stability}.
\end{proof}


\begin{corollary}[Exponential nonlinear feedback in exact conformal jets]
\label{cor:conformal-feedback-bound}
Fix the preceding background and let $f$ be a conformal map from the
disk to $D$, with $f(0)=0$. Fix $0<\sigma<1/2$, put
$q=\lfloor\sigma m\rfloor$, and let
\[
 E=\sum_{j=m-q+1}^m\ell_jU^j,\qquad V=\sum_j|\ell_j|.
\]
For large $m$, the difference between the perturbed corner
$e_0^*((G+E)^2-(1+f(u))U)^{-1}e_m$ and its affine approximation in
$E$ has every $u$-Taylor coefficient bounded by
\[
 \frac{C(m+1)^8 r^{-d}V^2}{1-Cr^{-d}V},\qquad
 d=m+2-2q-2p,
\]
provided $Cr^{-d}V<1$. Here $p$ is the fixed lower bandwidth of $G$
and $r>1$ is fixed. In particular $q=\lfloor3m/8\rfloor$ permits
more unknowns than the $m/4$ target jets while preserving an
exponentially small nonlinear remainder.
\end{corollary}
\begin{proof}
For large $m$, $E^2=0$. The support and gap argument in
Lemma~\ref{lem:corner-feedback} applies with its $h$ replaced by $q$.
Use $\|W_m\|\le C(m+1)^4$, $\|D_rW_mD_r^{-1}\|\le C$, and
$\|GE+EG\|\le CV$, uniformly on $\overline D$.
This gives the displayed supremum bound for the nonlinear remainder
throughout $D$. Composition with $f$ and Cauchy's estimate on circles
with radii increasing to one give the same bound for every coefficient,
without an exponentially growing coefficient-extraction factor.
\end{proof}


\begin{corollary}[Uniform replacement of the finite Jacobian on the boundary]
\label{cor:boundary-jacobian-replacement}
Fix $0<\sigma<1$ and $q=\lfloor\sigma m\rfloor$ in the preceding
setting. For $0\le d<q$, let $J_{m,d}(t)$ be the exact corner
derivative with respect to $g_{-(m-d)}$. Then
\[
 \sup_{t\in\overline D}
 \left|J_{m,d}(t)+2[s^d]R_t(s)X_t(s)\right|
 \le C(m+1)^3r^{-(m+1-q-2p)}.
\]
After composition with a fixed conformal map $f:\mathbb D\to D$,
the same bound holds for every Taylor coefficient of the difference.
In particular, no exponentially increasing coefficient-extraction
factor is needed when replacing the finite Jacobian in conformal jets.
\end{corollary}
\begin{proof}
Decompose the unsquared factor as
$p_w(z)=a_t(z^2)+wzb_t(z^2)$. Its coefficient functions, like those
of $A_t,B_t$, are holomorphic through $\overline D$.
Rationalizing the first-row generating function gives
\[
 R_t(s)=\frac{\zeta a_t(\zeta)+\zeta^2b_t(\zeta)g(\zeta)}
                 {h(\zeta)-t},\qquad
 X_t(s)=\frac{\zeta a_t(\zeta)g(\zeta)
                         +\zeta(1+t)b_t(\zeta)}{h(\zeta)-t},
 \quad \zeta=s^{-1}.
\]
Both are holomorphic in the disk, including at $s=0$ after
cancellation. The collar distance bound in Proposition~\ref{prop:boundary-polynomial-stability}
shows that each has modulus at most $C/(1-|s|)$, uniformly on
$\overline D$. Their coefficients are consequently bounded by $C(j+1)$.

Let $G_\infty=T(g)$. The formal infinite first-row equations, whose
sums in each column are finite because the positive bandwidth is $p$,
give $X_t=R_tG_\infty$ and
$R_t(G_\infty^2-(1+t)U_\infty)=e_0^*$.
These identities also hold at $t=-1$ by the displayed analytic formulas.
Truncate the row $R_t$ at length $L=m+1$. Its residual for the
finite pencil is supported in the last $2p$ columns and is bounded
by $CL$: all omitted indices are at most $L+2p$, and the fixed
Wiener norm of $g$ bounds the convolutions. The weighted pencil
inverse bound therefore gives
\[
 |(e_0^*W_m(t))_j-[s^j]R_t(s)|
       \le CLr^{j-L+2p}.
\]
For $j<q$ multiplication by $G$ gives the same estimate for the
finite $X$ row and $[s^j]X_t$, using $q+p<L$ and the finite
weighted Wiener norm of $g(r\,\cdot)$. In the coefficient of
degree $d<q$ of their product, summing at most $L$ terms now gives
the stated $CL^3r^{q-L+2p}$ bound. The exact derivative identity
\eqref{eq:odd-pencil-jacobian} supplies the factor $-2$.
Composition with $f$ gives a bounded analytic difference on the unit
disk; Cauchy's estimate with radii tending to one proves the final claim.
\end{proof}


\begin{proposition}[Exact limiting Jacobian as a boundary integral]
\label{prop:boundary-cauchy-jacobian}
For a polynomial $V(s)=\sum_{d\ge0}v_ds^d$, define
\[
 \mathcal J V(t)=-2\sum_{d\ge0}v_d[s^d]R_t(s)X_t(s),\qquad t\in D.
\]
Put
\[
 N_t(\zeta)=\zeta\left(A_t(\zeta)g(\zeta)
           +\tfrac12 B_t(\zeta)(h(\zeta)+t+2)\right).
\]
Then
\begin{equation}\label{eq:boundary-cauchy-jacobian}
 \mathcal J V(t)=-\frac2{2\pi i}\int_\T
          \frac{N_t(\zeta)V(\zeta)}{(h(\zeta)-t)^2}\,d\zeta.
\end{equation}
Let $\mathcal C_\Gamma f(t)=(2\pi i)^{-1}
\int_\Gamma f(x)/(x-t)\,dx$, and set
\[
 q(\zeta)=\frac{N_{h(\zeta)}(\zeta)}{h'(\zeta)},\qquad
 L_t(\zeta)=\frac{N_t(\zeta)-N_{h(\zeta)}(\zeta)}{t-h(\zeta)}.
\]
The latter quotient is removable at $t=h(\zeta)$. Formula
\eqref{eq:boundary-cauchy-jacobian} becomes
\[
 \mathcal J V=-2\,\partial_t\mathcal C_\Gamma
       [(qV)\circ h^{-1}]
   +2\,\mathcal C_\Gamma[(L_tV/h')\circ h^{-1}].
\]
Its leading weight has exactly the square-root endpoint zero:
\[
 q(\zeta)=\frac{\zeta g(\zeta)}{h'(\zeta)}
          p_{\sqrt\zeta\,g(\zeta)}(\sqrt\zeta)^2,
 \qquad
 \lim_{\zeta\to-1,\ \zeta\in\T}\frac{q(\zeta)}{g_0(\zeta)}
       =-p_0(i)^2\ne0.
\]
The expression is independent of the choice of $\sqrt\zeta$.
\end{proposition}
\begin{proof}
The odd part of $F_w(z)^2$ is $2wzR_t(z^2)X_t(z^2)$.
Use \eqref{eq:outside-square-parity} in $F_w^2$, put $s=z^2$
and $\zeta=s^{-1}$, and combine the two denominators. Direct algebra gives
\[
 R_t(\zeta^{-1})X_t(\zeta^{-1})
 =\frac{\zeta^2\{A_t(\zeta)g(\zeta)
        +\tfrac12 B_t(\zeta)(h(\zeta)+t+2)\}}
        {(h(\zeta)-t)^2}.
\]
Coefficient extraction proves the integral identity, first on a nearby
exterior circle and then on $\T$ by continuity. Splitting
$N_t=N_h+(t-h)L_t$ gives the two Cauchy terms with the displayed signs.
At $t=h(\zeta)$ take $w=\sqrt\zeta\,g(\zeta)$ in
\eqref{eq:outside-square-parity}; its right side is
$A_h+\zeta gB_h$. This proves the formula for $q$.
Finally $g/g_0\to1$, $h'\to1$ and $p_0$ is even with all its zeros
outside $\rho>1$. These facts prove the endpoint limit and its
nonvanishing.
\end{proof}


\section{Correction norms and endpoint restoration}
The base triangular Jacobian, indexed by reversed correction degree,
and its polynomial inverse are
\[
 (J^0v)_k=-2(k+1)\sum_{d\ge k}\binom{1/2}{d-k}v_d,
\]
\begin{equation}\label{eq:base-preconditioner}
 v_d=-\frac12\sum_{k=d}^{q-1}\binom{-1/2}{k-d}
                         \frac{x_k}{k+1},\qquad x=J^0v.
\end{equation}

\begin{lemma}[Localized invisible restoration and the squared curve]
\label{lem:localized-restoration}
Let $m\ge2$, let $L(s)=\sum_{j=1}^m\ell_js^{-j}$, and put
\[
 V=\sum_j|\ell_j|,\qquad
 T=\sum_{j\in\mathbb Z}|d_j-d_{j-1}|,
 \quad d_j=j(-1)^j\ell_j,
\]
where coefficients outside $1\le j\le m$ are zero. Define
\[
 K_m(s)=(-s)^{-m-1}\left(\frac{1-s^{-1}}2\right)^m,
 \qquad D(s)=L(s)-L(-1)K_m(s).
\]
Then $D(-1)=0$. The added restoring packet has coefficient
$\ell^1$ norm $|L(-1)|\le V$, and its frequencies in $zD(z^2)$
are invisible to $T_{2m+1}$. For $g_0(s)=\sqrt{1+s^{-1}}$,
\begin{equation}\label{eq:localized-weighted-c1}
 \|g_0D\|_{C^1(\T)}
 \le C\left(m^{3/4}V+\sqrt{mVT}\right).
\end{equation}
If $P$ is fixed and $C^1$ with $P(-1)=0$, then
\[
 \|s(g_0+P+D)^2-s(g_0+P)^2\|_{C^1}
 \le C_P\left(m^{3/4}V+\sqrt{mVT}+T+\sqrt m V+mV^2\right).
\]
All $C^1$ norms refer to the angular variable.
\end{lemma}
\begin{proof}
The binomial expansion has alternating coefficients of total absolute
sum one. Its support is $-m-1,\ldots,-2m-1$ in $s$, proving the
cost and invisibility statements.
Write $s=-e^{i\theta}$ and let $d$ be the distance of $\theta$ to
$2\pi\mathbb Z$. The square-root factor satisfies
$|g_0|\le C\sqrt d$ and $|\partial_\theta g_0|\le C/\sqrt d$.
Summation by parts and the trivial derivative bound give
\[
 |\partial_\theta L|\le C\min(mV,T/d),\qquad
 |L(s)-L(-1)|\le CV\min(md,1).
\]
The first bound follows by multiplying the derivative Fourier series
by $1-e^{-i\theta}$ and taking the absolute coefficient sum.

The packet has $|K_m|=|\cos(\theta/2)|^m$. Differentiating its explicit
formula, and using the Gaussian bound near zero and exponential decay
away from zero, gives
\[
 \sup\sqrt d\,|\partial_\theta K_m|\le Cm^{3/4},\qquad
 \sup d\,|\partial_\theta K_m|\le C\sqrt m,
 \qquad |1-K_m|\le C\min(md,1).
\]
Hence $|D|\le CV\min(md,1)$, so
$\|g_0'D\|_\infty\le C\sqrt m V$. Also
\[
 \sqrt d\,|L'|\le C\min(mV\sqrt d,T/\sqrt d)
                         \le C\sqrt{mVT}.
\]
These estimates prove \eqref{eq:localized-weighted-c1}.
At the endpoint the product is continuously differentiable, with
derivative zero, because $D(-1)=0$ and $D$ is differentiable.

Since $|P(s)|\le C_Pd$, the same bounds give
$\|PD\|_{C^1}\le C_P(V+T+\sqrt m V)$.
Finally $\|D\|_\infty\le2V$ and $\|D'\|_\infty\le CmV$, so
$\|D^2\|_{C^1}\le CmV^2$. Expanding the two squares proves the claim.
\end{proof}


\begin{lemma}[Coefficient and geometry bounds from the base preconditioner]
\label{lem:preconditioned-restoration}
Let $q\le m/2$, let $x$ have degree less than $q$, and define $v$
by \eqref{eq:base-preconditioner}. Put
\[
 L(s)=s^{-m}\sum_{d=0}^{q-1}v_ds^d,\quad
 D(s)=L(s)-L(-1)K_m(s),\quad
 R_1=\sum_{k<q}\frac{|x_k|}{k+1},\quad
 R_{1/2}=\sum_{k<q}\frac{|x_k|}{\sqrt{k+1}},
\]
where $K_m$ is the localized invisible packet from
Lemma~\ref{lem:localized-restoration}. Then
\[
 \|L\|_{\ell^1}\le C R_{1/2},\qquad
 \operatorname{Var}\{(-1)^dv_d\}_{d\in\mathbb Z}\le R_1,
 \qquad T_m\le C(mR_1+R_{1/2}).
\]
For $\gamma\ge0$ write
$\|a\|_{\mathcal W^\gamma}=\sum_j(1+|j|)^\gamma|a_j|$.
For every $0\le\alpha<1/2$,
\begin{equation}\label{eq:preconditioned-weighted-restoration}
 \|g_0D\|_{\mathcal W^{1+\alpha}}
 \le C_\alpha\bigl(m^{1+\alpha}R_1
                    +m^{3/4+\alpha}R_{1/2}\bigr).
\end{equation}
If $P$ is any fixed Laurent polynomial with $P(-1)=0$, then
\[
 \|s(g_0+P+D)^2-s(g_0+P)^2\|_{\mathcal W^{1+\alpha}}
 \le C_{P,\alpha}\bigl(m^{1+\alpha}R_1
       +m^{3/4+\alpha}R_{1/2}+m^{1+\alpha}R_{1/2}^2\bigr).
\]
\end{lemma}
\begin{proof}
Set $A_l=\binom{2l}{l}/4^l$. For one input coefficient $x_k$, the
modulated $v$ sequence is a multiple of the reversed finite sequence
$A_0,\ldots,A_k$. Its zero-padded variation is twice its maximum.
The factor $1/(2(k+1))$ in \eqref{eq:base-preconditioner} therefore
gives the stated variation bound. The identity
$\sum_{l=0}^kA_l=(2k+1)A_k=O(\sqrt{k+1})$ gives the coefficient
norm bound. Multiplication by the physical frequency $m-d$ gives
the bound for $T_m$.

A useful exact identity is, for $j\ge1$,
\[
 \sum_{d=0}^k\binom{1/2}{d+j}\binom{-1/2}{k-d}
  =\frac{k+1/2}{k+j}\binom{-1/2}{k}\binom{-1/2}{j-1}.
\]
It follows by differentiating
$(1+z)^{1/2}\sum_{l=0}^k\binom{-1/2}{l}z^l$.
The absolute values of these negative Laurent coefficients sum to
one: their alternating signs are all opposite to the surviving
nonnegative coefficient when evaluated at $s=-1$. Equivalently,
apply monotone convergence to the corresponding positive series
as its argument increases to one. Including the factor
$1/(2(k+1))$ shows that the full Wiener norm of $g_0$ times this
basis correction is $1/(k+1)$. Summing gives
\[
 \|g_0L\|_{\mathcal W^0}\le R_1.
\]

To estimate the restoring packet, modulate the variable to
$u=-s^{-1}$. Its binomial factor is $b_m(u)=((1+u)/2)^m$.
Unimodality gives $\|(1-u)b_m\|_{\ell^1}\le C m^{-1/2}$.
Write $(1-u)^{1/2}=1-\sum_{l\ge1}c_lu^l$, where
$c_l\ge0$, $\sum c_l=1$ and $c_l\le C l^{-3/2}$.
Then
\[
 \|(1-u)^{1/2}b_m\|_{\ell^1}
 \le\sum_{l\ge1}c_l\min(2,Clm^{-1/2})
 \le C m^{-1/4}.
\]
On frequencies of magnitude at most $4m+4$, these two Wiener
bounds, $|L(-1)|\le C R_{1/2}$, and multiplication by the largest
weight prove \eqref{eq:preconditioned-weighted-restoration}.

For completeness, the infinite tail is controlled by the endpoint
cancellation, not by a divergent first moment of $g_0$. In the
modulated variable $D$ is a polynomial supported at degrees at most
$2m+1$, with coefficient sum zero and coefficient norm at most
$CR_{1/2}$. For $n>4m+4$ the difference of two square-root kernel
coefficients whose indices differ by $j\le2m+1$ is bounded by
$Cj n^{-5/2}$, by the binomial recurrence. Thus the product's
$n$th coefficient is at most $CmR_{1/2}n^{-5/2}$. Its weighted tail
is bounded by $C_\alpha m^{1/2+\alpha}R_{1/2}$, which is smaller
than the restoring term already retained. This proves the first
weighted estimate.

Finally factor $P=(1+s)Q$ with $Q$ a fixed Laurent polynomial.
The variation bound gives $\|(1+s)L\|_{\ell^1}\le R_1$.
The binomial-packet difference bound gives
$\|(1+s)K_m\|_{\ell^1}\le Cm^{-1/2}$.
These are finite polynomials of frequency $O(m)$, so their
$\mathcal W^{1+\alpha}$ norms give the required bound for $PD$.
The term $D^2$ has coefficient norm at most $CR_{1/2}^2$ and
frequency $O(m)$. Expanding the squared curve proves the last assertion.
\end{proof}


\begin{corollary}[A sufficient negative-Sobolev correction size]
\label{cor:negative-sobolev-correction}
Fix $-1/2<s<0$ and $q\le m/2$. If
$\|x\|_{H^s}\le C_Pm^{s-1}$, the preceding correction satisfies
\[
 \|L\|_{\ell^1}=O_P(m^{-1}),\qquad T_m=O_P(m^s).
\]
For every $0\le\alpha<\min(-s,1/4)$ its restored squared-curve
change tends to zero in $\mathcal W^{1+\alpha}$, with bound
\[
 C_{P,s,\alpha}
       \bigl(m^{s+\alpha}+m^{\alpha-1/4}+m^{\alpha-1}\bigr).
\]
Consequently the construction can preserve a small $C^{1,\alpha}$
neighborhood, not just a $C^1$ neighborhood, if this norm estimate
for the preconditioned correction is established.
\end{corollary}
\begin{proof}
Weighted Cauchy--Schwarz gives
$R_1\le C_s\|x\|_{H^s}$ and
$R_{1/2}\le C_s q^{-s}\|x\|_{H^s}$.
The former uses $s>-1/2$, the latter $s<0$.
Since $q\le m/2$, these give $R_1=O_P(m^{s-1})$ and
$R_{1/2}=O_P(m^{-1})$. Substitute in the preceding lemma.
Absolute weighted Fourier summability implies the stated
$C^{1,\alpha}$ control.
\end{proof}


\begin{lemma}[The restored perturbation divided by the endpoint square root]
\label{lem:normalized-endpoint-restoration}
For the localized restoration of any negative correction $L$ supported
in frequencies of magnitude at most $m$, put $V=\|L\|_{\ell^1}$.
For every $0\le r<1$, in the two-sided circle Sobolev norm,
\[
 \|D/g_0\|_{H^r(\T)}
       \le C_r m^r V\sqrt{\log(m+1)}.
\]
In particular, corrections with $V=O_P(m^{-1})$ change $P/g_0$
by an arbitrarily small amount in $H^r(\T)$ as $m\to\infty$.
For $1/2<r<1$ this also controls the supremum norm.
\end{lemma}
\begin{proof}
The restored polynomial has $\|D\|_{\ell^1}\le2V$, frequencies
of magnitude at most $2m+1$, and $D(-1)=0$. At angular distance $d$
from $-1$, these facts give
$|D|\le CV\min(md,1)$, while $|g_0|^2$ is comparable to $d$.
Integration therefore gives
$\|D/g_0\|_2\le CV\sqrt{\log(m+1)}$.
After modulation to $u=-s^{-1}$, the coefficients of $1/g_0$ are
$A_n=\binom{2n}{n}/4^n$. Their first differences are $O(n^{-3/2})$.
The zero sum of the modulated coefficients of $D$ thus bounds the
$n$th product coefficient by $CmVn^{-3/2}$ for $n>4m+4$.
The low-frequency part of the $H^r$ norm is at most $Cm^r$ times
the $L^2$ norm. The squared high-frequency part is bounded by
$C m^2V^2\sum_{n>4m}n^{2r-3}\le C_rm^{2r}V^2$ for $r<1$.
These estimates prove the assertion. The last statement is the
coefficient Cauchy--Schwarz embedding for $r>1/2$.
\end{proof}


\section{The geometric operator estimate}
The following estimates control the actual Jacobian, including its
exterior-zero factors, uniformly as the degree of the background grows.

\begin{lemma}[The square-root extension of the analytic projection]
\label{lem:endpoint-extension}
Let $g_0(s)=\sqrt{1+s^{-1}}$ and let $y$ be a polynomial. Define
$V=T(g_0)^{-1}y$, using the triangular polynomial inverse, and
$Ey=g_0V$. For every $0<r<1$,
\[
 \|Ey\|_{H^r(\T)}\le C_r\|y\|_{H^r},\qquad \Pi_+Ey=y.
\]
For $r>1/2$, the same inverse also satisfies
$\|V\|_{L^2}\le C_r\|y\|_{H^r}$.
These operators therefore extend by density to the indicated spaces.
For $r>1/2$ the continuous representative of $Ey$ vanishes at $-1$.
\end{lemma}
\begin{proof}
Write $A_l=\binom{2l}{l}/4^l$. The exact partial-convolution identity
in Lemma~\ref{lem:preconditioned-restoration} gives, up to signs, the
matrix of the negative part of $E$:
\[
 K_{j,k}=\frac{(k+1/2)A_kA_{j-1}}{k+j},
 \qquad j\ge1,\quad k\ge0.
\]
Since $A_l\le C(l+1)^{-1/2}$, the matrix after the Sobolev weights
is bounded entrywise by
\[
 C\,\frac{j^{r-1/2}(k+1)^{1/2-r}}{j+k}.
\]
Apply Schur's test with row weights $j^{-1/2}$ and column weights
$(k+1)^{-1/2}$. Its two required sums are bounded by
\[
 Cj^r\sum_{k\ge0}\frac{(k+1)^{-r}}{j+k},
 \qquad
 C(k+1)^{1-r}\sum_{j\ge1}\frac{j^{r-1}}{j+k},
\]
respectively. Splitting each sum at $j=k+1$, or comparing with
the corresponding integral, bounds both uniformly for $0<r<1$.
This proves the first assertion.

The coefficient matrix of $T(g_0)^{-1}:H^r\to L^2$ has squared
Hilbert--Schmidt norm
\[
 \sum_{k\ge0}(k+1)^{-2r}\sum_{d=0}^k A_{k-d}^2
 \le C\sum_{k\ge0}(k+1)^{-2r}\log(k+2)<\infty
\]
when $r>1/2$. For polynomial inputs $Ey=g_0V$ vanishes at $-1$.
Point evaluation is continuous on $H^r(\T)$ for $r>1/2$,
so the last assertion follows by density.
\end{proof}


\begin{lemma}[Small Cauchy-projection perturbations in the circle parameter]
\label{lem:wiener-cauchy-projections}
Let $h(s)=s+k(s)$, with
$k\in\mathcal W^{1+\alpha}$ for some $\alpha>0$, and put
$\delta=\sum_n |n|\,|k_n|$. If $\delta$ is sufficiently small,
$h(\T)$ is a positively oriented Jordan curve. Let $C_h$ denote
the interior Cauchy projection pulled back by $h$. For $0\le r\le1$,
\[
 \|C_h-\Pi_+\|_{H^r(\T)\to H^r(\T)}\le C_r\delta .
\]
Write $k=k_-+k_+$, where $k_-$ has nonpositive powers and $k_+$
has strictly positive powers. If, in addition,
$\|k_+\|_{\mathcal W^{r+1}}$ is sufficiently small, then
\[
 \|\Pi_+C_h\Pi_-\|_{L^2\to H^r}
       \le C_r\|k_+\|_{\mathcal W^{r+1}}.
\]
Here $\Pi_-=I-\Pi_+$; the constants do not involve higher derivatives
of $k_-$.
\end{lemma}
\begin{proof}
The divided difference
\[
 H(\zeta,\eta)=\frac{h(\zeta)-h(\eta)}{\zeta-\eta}
\]
belongs to the two-variable Wiener algebra and
$\|H-1\|_{\mathcal W(\T^2)}\le\delta$. Indeed the divided
difference of a monomial of degree $n$ has $|n|$ terms of
coefficient modulus one, for either sign of $n$. Thus
$|h(\zeta)-h(\eta)|\ge(1-\delta)|\zeta-\eta|$.
The same statement along the homotopy $s+tk(s)$ proves the
orientation assertion. The constant term of $k$ causes only a
translation.

Set $A(\zeta,\eta)=h'(\zeta)/H(\zeta,\eta)$, where $h'$ is the
derivative with respect to the circle variable. The Wiener
Neumann series gives
\[
 \|A-1\|_{\mathcal W(\T^2)}
       \le\frac{2\delta}{1-\delta}.
\]
If $A=\sum a_{j,l}\zeta^j\eta^l$, the Cauchy boundary formula is
the norm-convergent operator series
\[
 C_h=\sum_{j,l}a_{j,l}M_{s^l}\Pi_+M_{s^j}.
\]
The half-jump terms agree because $A(\eta,\eta)=1$.
This identity follows first with radial regularization of the
circle Cauchy kernel and then by its $L^2$ limit; the
$C^{1,\alpha}$ boundary has the usual Cauchy boundary values.
It proves the asserted small bound on $L^2$.
The operator is a projection. This follows from the Cauchy formula
for smooth curves and then from Fourier approximation of $h$:
the displayed Wiener representation gives operator-norm convergence
on $L^2$, so the identity $C_h^2=C_h$ passes to the limit.

For smooth $v$, integration by parts in the Cauchy integral gives
\[
 \partial_\theta(C_hv)
  =h_\theta\,C_h(v_\theta/h_\theta).
\]
The multipliers $h_\theta$ and $h_\theta^{-1}$ are uniformly
bounded on $L^2$ and are within $C\delta$ there of $is$ and
$(is)^{-1}$. Subtracting the analogous identity for $\Pi_+$ and
using the $L^2$ estimate therefore gives the small $H^1$ bound.
The formula and bound for the stated regularity follow by Fourier
approximation of $h$ and density of smooth $v$. Interpolation
between $L^2$ and $H^1$ proves the estimate for $0<r<1$.

For the last claim use the two-variable algebra with coefficient
weight $(1+\max(j,0)+\max(l,0))^r$. The divided differences of
$k_-$ have only nonpositive powers in each variable, and their
weighted norm is at most $\delta$. Those of $k_+$ have weighted
norm at most $C_r\|k_+\|_{\mathcal W^{r+1}}$.
Neumann expansion consequently splits $A=A_-+A_{\rm rem}$, where
$A_-$ has only nonpositive powers in both variables and
\[
 \|A_{\rm rem}\|_{\mathcal W_{\rm pos}^r(\T^2)}
      \le C_r\|k_+\|_{\mathcal W^{r+1}}.
\]
Every term of $\Pi_+C_h\Pi_-$ coming from $A_-$ vanishes. A term
$\Pi_+M_{s^l}\Pi_+M_{s^j}\Pi_-$ can be nonzero only for $j\ge1$
and $j+l\ge1$. It is a partial shift from negative input
frequencies to output frequencies between $0$ and $j+l-1$.
Its $L^2\to H^r$ norm is therefore at most
$(1+\max(j,0)+\max(l,0))^r$. Summing the remainder coefficients
proves the final estimate.
\end{proof}


\begin{lemma}[An endpoint Bezout kernel in the Wiener algebra]
\label{lem:endpoint-bezout}
Let $b(x)=\sqrt{1+x}$ and let $D$ be a polynomial with $D(-1)=0$.
For every $\alpha>0$,
\[
 \left\|\frac{b(y)D(x)-b(x)D(y)}{y-x}\right\|_{\mathcal W(\T^2)}
 \le C_\alpha\|bD\|_{\mathcal W^{1+\alpha}}.
\]
The assertion is used only when its right side is finite.
\end{lemma}
\begin{proof}
Put $v=bD$ and $q=v/(1+x)$. The kernel on the left is
$b(x)b(y)(q(x)-q(y))/(y-x)$. Since $v(-1)=0$, write
\[
 v(x)=\sum_{n\ge1}v_n\bigl(x^n-(-1)^n\bigr).
\]
For one summand $x^n-(-1)^n$, its quotient by $1+x$ is an
alternating geometric polynomial. After replacing $x,y$ by $-x,-y$,
its kernel is, up to sign,
\[
 \sqrt{1-x}\sqrt{1-y}\,T_{n-2}(x,y),\qquad
 T_N=\sum_{\substack{i,j\ge0\\i+j\le N}}x^iy^j.
\]
The case $n=1$ has zero kernel.

Write $\sqrt{1-x}=\sum_{\ell\ge1}c_\ell(1-x^\ell)$, where
$c_\ell>0$, $\sum c_\ell=1$, and
$\sum_{\ell\ge t}c_\ell=\binom{2t-2}{t-1}/4^{t-1}\le C/\sqrt t$.
Regard the coefficients of $T_{n-2}$ as a triangular indicator
on $\mathbb Z^2$. A horizontal translation by $\ell$ changes its
$\ell^1$ norm by at most $Cn\min(\ell,n)$, and the analogous
vertical bound holds. Consequently
\[
 \|(I-S_x^\ell)(I-S_y^k)T_{n-2}\|_{\ell^1}
       \le Cn\min(\ell,k,n).
\]
Summing against $c_\ell c_k$ gives
\[
 \|\sqrt{1-x}\sqrt{1-y}\,T_{n-2}\|_{\mathcal W}
 \le Cn\sum_{t=1}^n\left(\sum_{\ell\ge t}c_\ell\right)^2
 \le Cn\log(n+1).
\]
Finally $n\log(n+1)\le C_\alpha(1+n)^{1+\alpha}$.
The resulting absolutely convergent kernel series agrees with the
displayed divided difference in the open bidisk and hence on its
continuous boundary extension.
\end{proof}


\begin{lemma}[The unsquared divided difference]
\label{lem:unsquared-divided-difference}
Write $g=g_0+P_-+P_+$, where $P_-$ has strictly negative powers,
$P_+$ has strictly positive powers, and both vanish at $-1$.
For $a(z)=zg(z^2)$ and $a_0(z)=zg_0(z^2)$, set
\[
 F(z,\eta)=\frac{a(z)-a(\eta)}{a_0(z)-a_0(\eta)}
          =F_-(z,\eta)+F_+(z,\eta),
\]
where $F_--1$ and $F_+$ are the contributions of $P_-$ and $P_+$.
All quotients have continuous extensions at their removable boundary
points. For $\alpha>0$,
\[
 \|F_--1\|_{\mathcal W(\T^2)}
       \le C_\alpha\|g_0P_-\|_{\mathcal W^{1+\alpha}},
\]
and $F_--1$ has strictly negative powers of each variable.
For $K\ge0$, let $\mathcal W_{\mathrm{pos},z}^K$ have coefficient
weight $(1+\max(j,0))^K$ on the exponent $j$ of $z$, with no
weight on the exponent of $\eta$. Then
\[
 \|F_+\|_{\mathcal W_{\mathrm{pos},z}^K}
       \le C_K\|P_+\|_{\mathcal W^{K+2}}.
\]
\end{lemma}
\begin{proof}
Rationalizing the negative contribution separates it into
\[
 \frac{s g_0(s)P_-(s)-t g_0(t)P_-(t)}{s-t}
 +z\eta\,\frac{g_0(t)P_-(s)-g_0(s)P_-(t)}{s-t},
 \qquad s=z^2,\quad t=\eta^2.
\]
The first divided difference has Wiener norm bounded by
$C\|g_0P_-\|_{\mathcal W^1}$. For the second put
$x=s^{-1}$, $y=t^{-1}$ and $D(x)=P_-(x^{-1})$.
It becomes $z^{-1}\eta^{-1}$ times the kernel in
Lemma~\ref{lem:endpoint-bezout}. The claimed norm and Fourier
orientation follow. In particular the quotient extends at the
two distinct contact parameters where both base values are zero.

For the positive contribution write $P_+(s)=(1+s)Q(s)$ and
$r(z)=zQ(z^2)$, an odd polynomial. Put $A=a_0(z)$ and $B=a_0(\eta)$.
Since the change in $a$ is $A^2r(z)$, exact algebra gives
\[
 F_+=A r(z)+B r(\eta)
   +AB\,\frac{A+B}{z+\eta}\,
                   \frac{r(z)-r(\eta)}{z-\eta}.
\]
The two-variable kernel $AB(A+B)/(z+\eta)$ has finite Wiener
norm and has at most one positive power in each variable.
Indeed, with $x=z^{-1}$, $y=\eta^{-1}$,
$b_x=\sqrt{1+x^2}$ and $b_y=\sqrt{1+y^2}$, it equals
\[
 (xy)^{-1}\left((1+x^2)b_y
       +x(y-x)\frac{b_xb_y}{b_x+b_y}\right).
\]
The last fraction belongs to the bidisk Wiener algebra because
\[
 \frac{b_xb_y}{b_x+b_y}
   =\int_0^\infty b_xe^{-ub_x}b_ye^{-ub_y}\,du
\]
is a convergent Bochner integral there. To verify convergence,
the elementary subordination formula
\[
 e^{-u\sqrt{1+x^2}}
 =\frac{u}{2\sqrt\pi}\int_0^\infty
       v^{-3/2}e^{-u^2/(4v)}e^{-v(1+x^2)}\,dv
\]
gives Wiener norm at most one. Its positive scalar density has
mass one, and its $u$ derivative has total variation $C/u$,
by the substitution $v=u^2w$. Thus
$\|b_xe^{-ub_x}\|_{\mathcal W}\le C/u$; the bound $2$ for
small $u$ follows instead from $\|b_x\|_{\mathcal W}\le2$.
The product is integrable in $u$. The formula is first checked
in the open bidisk, where the square roots have positive real
part, and then extended by the Wiener convergence.

A monomial of degree $n$ in $r$ now contributes at most
$C_K(1+n)^{K+1}$ to the required norm: its divided difference
has $n$ terms, and the resulting positive $z$ degree is at
most $n+1$. Finally division by $1+s$, using $P_+(-1)=0$,
gives $\|Q\|_{\mathcal W^{K+1}}\le
C_K\|P_+\|_{\mathcal W^{K+2}}$. This proves the last assertion.
\end{proof}


\begin{lemma}[A quantitative conformal coordinate near the circle]
\label{lem:quantitative-conformal-coordinate}
For fixed $0<\alpha<1$, a sufficiently small
$\mathcal W^{1+\alpha}$ neighborhood of $h(s)=s$ admits normalized
Riemann maps $f:\mathbb D\to D$ with
\[
 \|f'-1\|_\infty\le C_\alpha\|h-s\|_{\mathcal W^{1+\alpha}}.
\]
The maps have uniformly bounded $C^{1,\alpha}$ norms on the closed
disk. The smallness threshold and bound follow from elementary
H\"older norm estimates and can be made effective.
For the real-symmetric curves used here, $f(-1)=-1$.
\end{lemma}
\begin{proof}
Smallness in the stated Fourier norm implies smallness in
$C^{1,\alpha}(\T)$. The angle of $h(e^{it})$ is a
$C^{1,\alpha}$ diffeomorphism close to $t$. Inverting it
represents the curve as $\rho(\theta)e^{i\theta}$, with
\[
 \|\ell\|_{C^{1,\alpha}}\le C_\alpha\|h-s\|_{\mathcal W^{1+\alpha}},
 \qquad \ell=\log\rho .
\]
This estimate follows directly by differentiating the angle and
the logarithm of $h(e^{it})e^{-it}$, and applying the product,
composition, and inverse-function estimates in $C^{1,\alpha}$.
All denominators stay uniformly away from zero.

Let $H$ be the periodic Hilbert transform, with zero mean and
the sign for which it supplies the imaginary part of an analytic
function from its real part. It is bounded on $C^\alpha$.
For completeness, subtract the value at the evaluation point
in its cotangent-kernel formula; splitting at the distance
between two evaluation points bounds the H\"older seminorm by
$C(\alpha^{-1}+(1-\alpha)^{-1})$ times the input seminorm.
It commutes with differentiation and has the corresponding
bound on $C^{1,\alpha}$.

Seek a real mean-zero angle correction $u$ satisfying
\[
 u=H[\ell(t+u(t))].
\]
Use the closed set $\|u\|_{C^{1,\alpha}}\le M$, where
$M=C_\alpha\|\ell\|_{C^{1,\alpha}}<1/4$, with a sufficiently
large fixed coefficient, but give this set the $C^\alpha$ metric.
The set is complete: a $C^\alpha$ limit with these uniform
derivative bounds is still $C^{1,\alpha}$ with the same bound,
by compactness of the derivatives and the fundamental theorem
of calculus. The Hilbert-transform and composition estimates
show that the displayed map preserves the set for small $\ell$.
For two inputs use
\[
 \ell(t+u)-\ell(t+v)
  =(u-v)\int_0^1\ell'(t+v+\lambda(u-v))\,d\lambda.
\]
The intervening angle maps are uniformly Lipschitz. H\"older
product and composition bounds therefore give Lipschitz
constant $C_\alpha\|\ell\|_{C^{1,\alpha}}<1$ in the
$C^\alpha$ metric. The contraction theorem supplies the
required small $C^{1,\alpha}$ solution.

Let $G$ be analytic in the disk with boundary real part
$\ell(t+u(t))$, boundary imaginary part $u(t)$, and real $G(0)$.
Then $f(z)=z\exp G(z)$ has the prescribed boundary curve,
$f(0)=0$, and $f'(0)>0$. The norm bounds for $G$ give the
displayed derivative estimate; all derivatives extend
continuously to the boundary. Poisson extension preserves the
$C^\alpha$ class on the closed disk: subtract the value at the
nearest boundary point and split the Poisson integral at angular
distance $1-|z|$ to bound the radial error by
$C_\alpha(1-|z|)^\alpha$; angular differences follow by translation.
Apply this to the boundary derivative of $G$.
If the neighborhood is small,
$\|f'-1\|_\infty<1/4$, so integration along disk chords proves
univalence and uniform bi-Lipschitz bounds. The boundary
homeomorphism and the argument principle identify its image
with $D$. In the symmetric case $\ell$ is even. Uniqueness
makes $u$ odd, hence $u(\pi)=0$, and $\rho(\pi)=1$ gives
$f(-1)=-1$.
\end{proof}


\begin{proposition}[Uniform small far factors and spectral derivatives]
\label{prop:small-far-factors}
Under the hypotheses of Lemma~\ref{lem:unsquared-divided-difference}, suppose that
$\|g_0P_-\|_{\mathcal W^{1+\alpha}}$ and
$\|P_+\|_{\mathcal W^{K+3}}$ are sufficiently small.
Assume also the preceding squared-curve neighborhood, so that
the normalized $f$ and its inverse have uniform bi-Lipschitz bounds.
For the actual far factors
\[
 p_w(z)^2=A_t(z^2)+wzB_t(z^2),\qquad t=w^2-1,
\]
one has, for $K\ge0$,
\[
 \sum_{j\ge0}(1+j)^K\left(
  \|[s^j](A_t(s)-1)\|_{H^\infty(D)}
       +\|[s^j]B_t(s)\|_{H^\infty(D)}\right)
 \le C_K\|P_+\|_{\mathcal W^{K+3}}.
\]
The constants are uniform on a sufficiently small fixed neighborhood.
\end{proposition}
\begin{proof}
For a boundary value $w=a(\eta)$ the positive Wiener factor is
\[
 \log p_{a(\eta)}(z)=\Pi_{>0,z}\log F(z,\eta).
\]
This can be checked on a sufficiently thin exterior circle:
$(a_0(z)-a_0(\eta))/z$ has a zero-free exterior logarithm
with no positive Laurent part, and the positive part of the
remaining logarithm is the logarithm of the normalized far factor.
The previously proved collar separates its zeros from that circle.
Taking its boundary limit gives the displayed identity in the
two-variable Wiener algebra.

The function $F_-$ is close to one and coanalytic in both
variables. Its logarithm contributes nothing to $\Pi_{>0,z}$.
Neumann series in the submultiplicatively weighted algebra
$\mathcal W_{\mathrm{pos},z}^K$, followed by exponentiation, give
\[
 \|p_{a(\eta)}(z)-1\|_{\mathcal W_{\mathrm{pos},z}^K}
       \le C_K\|P_+\|_{\mathcal W^{K+2}}.
\]
In particular this bounds the sum of the weighted suprema of
the individual positive-$z$ coefficients on the spectral boundary.
Their holomorphic continuation in $w$, established in the collar
lemma, and the maximum principle give the same bound on the two
closed lobes $\{w:w^2-1\in\overline D\}$.

For the odd coefficients we also need a spectral derivative bound.
Differentiating the factor on an exterior circle gives
\[
 \partial_w\log p_w(z)=-\Pi_{>0,z}\frac1{a(z)-w}.
\]
At $w=a(\eta)$ the reciprocal is
\[
 \frac1{a_0(z)-a_0(\eta)}\,F(z,\eta)^{-1}.
\]
The first factor has only negative powers of $z$; every one of
its coefficients has Wiener norm at most $2$ as a function of
$\eta$. To see this, expand
$(a_0(z)+a_0(\eta))/(z^2-\eta^2)$ geometrically in $z^{-2}$.
The odd coefficients are partial sums of square-root binomial
coefficients times powers of $\eta$, and the even coefficients
are $a_0(\eta)$ times a power of $\eta$.

Subtract $F_-^{-1}$ from $F^{-1}$: its product with the base
reciprocal is coanalytic and disappears under the positive
projection. For any remaining series $E$ the coefficient count
just described gives
\[
 \left\|\Pi_{>0,z}
       \frac{E(z,\eta)}{a_0(z)-a_0(\eta)}
       \right\|_{\mathcal W_{\mathrm{pos},z}^K}
 \le C_K\|E\|_{\mathcal W_{\mathrm{pos},z}^{K+1}}.
\]
Only positive input degrees $j$ contribute, and at most $j$
base coefficients can contribute for each such degree. The
series is absolutely convergent in the indicated norm.
For the fixed Laurent-polynomial background it also agrees
with the exterior-circle projection: the positive Laurent
coefficients decay exponentially on its fixed collar, so the
radial limit can be taken term by term. Apply the estimate to
$E=F^{-1}-F_-^{-1}$ and use the preceding Neumann bounds. It
follows that
\[
 \|\partial_w p_{a(\eta)}(z)\|_{\mathcal W_{\mathrm{pos},z}^K}
       \le C_K\|P_+\|_{\mathcal W^{K+3}}.
\]
The same holds for $\partial_w(p_w^2)$, and the maximum principle
again bounds the sum of coefficient suprema throughout both lobes.

For an odd coefficient $Q_j(w)=[z^j]p_w(z)^2$, parity gives
$Q_j(0)=0$. There is a path in the corresponding closed lobe
from $0$ to $w$ of length at most $C|w|$. Indeed, if
$w^2-1=f(u)$, use the path
$w(v)=\sqrt{1+f(-1+v(u+1))}$, $0\le v\le1$, with the
continuous choice of sign. Bi-Lipschitz bounds give
$|1+f(-1+v(u+1))|\asymp v|u+1|$, so integration of its
derivative bounds the length by $C\sqrt{|u+1|}\le C|w|$.
Consequently $|Q_j(w)/w|\le C\sup|Q'_j|$.
The even coefficient bound supplies $A_t-1$, and the odd
coefficient bound just proved supplies $B_t$. Their weighted
sums are exactly the asserted estimate.
\end{proof}


\begin{proposition}[The two-level isomorphism estimate]
\label{prop:two-level-isomorphism}
Fix $\alpha=1/8$ and split the finite real Laurent background as
$P=P_-+P_+$, with strictly negative and strictly positive powers,
respectively, and with $P_-(-1)=P_+(-1)=0$. There is a sufficiently
small $\gamma>0$ such that the bounds
\[
\begin{split}
 &\|h-s\|_{\mathcal W^{1+\alpha}}<\gamma,\qquad
   \|P_-\|_{\mathcal W^0}<\gamma,\\
 &\|P_-/g_0\|_{\mathcal W^0}<\gamma,\qquad
   \|g_0P_-\|_{\mathcal W^{1+\alpha}}<\gamma,\qquad
   \|P_+\|_{\mathcal W^4}<\gamma
\end{split}
\]
imply that the actual operator
$\mathcal A_gx=(\mathcal J_g(J^0)^{-1}x)\circ f$ extends to
compatible isomorphisms on $H^{-1/4}$ and $H^{-1/8}$.
Their norms and inverse norms are bounded independently of the
degree of $P$.
\end{proposition}
\begin{proof}
The proof applies at each of $r=3/4,7/8$, with $s=r-1$.
For a one-variable multiplier use the coefficient norm
\[
 \|M\|_{\mathcal W_{\rm pos}^r}
   =\sum_{n\in\mathbb Z}(1+\max(n,0))^r|M_n|.
\]
It is submultiplicative. The endpoint extension lemma gives
\begin{equation}\label{eq:one-sided-density-product}
 \|\Pi_+(M Ey)\|_{H^r}
       \le C_r\|M\|_{\mathcal W_{\rm pos}^r}\|y\|_{H^r},
 \qquad
 \|M Ey\|_{L^2}\le C_r\|M\|_{\mathcal W^0}\|y\|_{H^r}.
\end{equation}
Indeed a coanalytic multiplier cannot turn the negative part of
$Ey$ positive, and its backward shifts on analytic $H^r$ are
contractions. A positive shift by $j$ on two-sided $H^r$ costs
at most $C_r(1+j)^r$. This proves both assertions by summation.

Put $R=g/g_0$ and $m=R/h'$. The assumptions imply
\[
 \|m-1\|_{\mathcal W_{\rm pos}^r}\le C_r\gamma,\qquad
 \|(h-s)_+\|_{\mathcal W^{r+1}}\le C_r\gamma .
\]
Here are the Fourier details. The negative part
$1+P_-/g_0$ of $R$ is close to one in $\mathcal W^0$.
Writing $P_+=(1+s)Q$ gives $P_+/g_0=sg_0Q$ and
$\|Q\|_{\mathcal W^r}\le C_r\|P_+\|_{\mathcal W^{r+1}}$.
The part of $h-s$ without positive-background factors,
$2sg_0P_-+sP_-^2$, is nonpositive. Every other term is a
product containing $P_+$, so its positive weighted norm of
order $r+1$ is bounded by $C_r\|P_+\|_{\mathcal W^{r+2}}$,
using the bounded unweighted norms of the coanalytic factors.
Thus $h'-1$ is small in $\mathcal W_{\rm pos}^r$; Neumann
inversion proves the bound for $m$.

Consider the operator on analytic $H^r$
\[
 By=\Pi_+s^{-1}C_h[s m Ey].
\]
Its direct projection is $\Pi_+m Ey$, within $C_r\gamma$ of $y$
by \eqref{eq:one-sided-density-product}. The positive part of
the density $s m Ey$ is in $H^r$ with bounded norm. Apply
the small $H^r$ estimate for $C_h-\Pi_+$ to that part, and the
$L^2\to H^r$ estimate for $\Pi_+C_h\Pi_-$ to its negative part.
Lemma~\ref{lem:wiener-cauchy-projections} then gives
\[
 \|B-I\|_{H^r\to H^r}\le C_r\gamma .
\]
Consequently $B$ is an isomorphism for small $\gamma$.

We explain why this is the required bare Cauchy isomorphism,
rather than merely a projected estimate. Let
$\mathcal R_r=C_hH^r(\T)$. The small norm of $C_h-\Pi_+$ implies
that $\Pi_+:\mathcal R_r\to H^r$ is an isomorphism: on analytic
inputs $\Pi_+C_h$ is close to the identity, while a range vector
$v$ with $\Pi_+v=0$ satisfies $v=(C_h-\Pi_+)v$ and is zero.
Constants are fixed by $C_h$. Therefore
\[
 \pi(v)=\Pi_+s^{-1}v
\]
identifies $\mathcal R_r$ modulo constants with analytic $H^r$.
The same range argument on $L^2$ shows that an $L^2$ range
vector whose positive part is in $H^r$ actually lies in
$\mathcal R_r$. This applies to the densities just considered.

Composition by $h^{-1}\circ f$ identifies $\mathcal R_r$ with
analytic $H^r$, preserving constants, with uniform bounds in
both directions. The boundary maps are uniformly bi-Lipschitz.
For $0<r<1$ their action on $H^r(\T)$ follows from the
equivalent norm
\[
 \|v\|_{L^2}^2+
  \int_\T\int_\T
     \frac{|v(z)-v(\zeta)|^2}{|z-\zeta|^{1+2r}}
                       \,|dz|\,|d\zeta|,
\]
by change of variables. Equivalence with the Fourier norm
follows by expanding in Fourier modes. The range identification
then follows from the Cauchy formula, first for analytic
polynomials in the disk coordinate and then by density.

Differentiation in the physical spectral variable, followed by
composition with $f$, acts as $Z\mapsto Z'/f'$ in the disk.
Differentiation is an isomorphism from analytic $H^r$ modulo
constants onto $H^{r-1}$. Analytic $H^\infty$ multipliers are
bounded on $H^s$ for $s<0$: this follows from the equivalent
weighted Bergman norm
\[
 \int_{\mathbb D}|X(u)|^2(1-|u|^2)^{-1-2s}\,dA(u).
\]
The equivalence is immediate by integrating each monomial
and comparing its beta-integral weight with $(k+1)^{2s}$.
Both $f'$ and $1/f'$ are uniformly bounded here.

Define the bare operator by
\[
 \mathcal J_{\rm bare}V
  =-2\,\partial_t\mathcal C_\Gamma[(sgV/h')\circ h^{-1}].
\]
Since $y=\Pi_+g_0V$ satisfies $x_k=-2(k+1)y_k$ when
$x=J^0V$, the preceding isomorphisms factor
$(\mathcal J_{\rm bare}(J^0)^{-1}x)\circ f$
into invertible maps with uniform bounds. This proves the bare
isomorphism on $H^s$. The same proof of boundedness, without
requiring a multiplier near one, shows that the differentiated
Cauchy transform of a density $sM Ey$ has $H^s$ output norm
at most $C_r\|M\|_{\mathcal W_{\rm pos}^r}\|y\|_{H^r}$.

It remains to retain the actual exterior factors. Write
$A_t(s)=\sum A_j(t)s^j$ and $B_t(s)=\sum B_j(t)s^j$ in
Proposition~\ref{prop:boundary-cauchy-jacobian}.
The $A$ part of the actual Jacobian is
$\sum_j A_j(t)\mathcal J_{\rm bare}(s^jV)(t)$.
The shift costs at most $C_r(1+j)^r$ in the preceding density
bound. Its coefficient is an analytic multiplier on $H^s$
after composition with $f$, so it costs only its $H^\infty(D)$
norm.

For the $B$ part use the exact identities
$h+t+2=2(h+1)-(h-t)$ and $h+1=sg^2$.
The first resulting term is the differentiated Cauchy transform
with density $s^{j+2}g^2V/h'=s^{j+2}gm Ey$.
The general density bound applies with $M=s^{j+1}gm$, at
cost at most $C_r(1+j)^r$. The second term is the
undifferentiated Cauchy transform with density $s^{j+1}V/h'$.
It is bounded from $y\in H^r$ to $L^2$, by the second assertion
of Lemma~\ref{lem:endpoint-extension}, boundedness of $C_h$ on
$L^2$, and the bounded multiplier $1/h'$. The boundary
composition is bounded on $L^2$, which embeds in $H^s$.
This bound does not grow with $j$.

Proposition~\ref{prop:small-far-factors}, used with $K=r\le1$,
now proves
\[
 \|\mathcal A_g-\mathcal A_{{\rm bare},g}\|_{H^s\to H^s}
       \le C_r\|P_+\|_{\mathcal W^4}\le C_r\gamma .
\]
Neumann inversion, after decreasing the single fixed $\gamma$,
gives the asserted actual isomorphisms for both chosen exponents.
All formulas agree on polynomial inputs and extend compatibly.
The inverse on the stronger space agrees with the weaker inverse
by uniqueness in the weaker space.
\end{proof}


\section{The forcing and growing finite systems}
\begin{proposition}[The scattering generating function]
\label{prop:fixed-background-forcing}
With $\zeta=s^{-1}$ set
\[
 S_t(s)=\frac{\zeta A_t(\zeta)+\zeta^2B_t(\zeta)g(\zeta)}
                  {h(\zeta)-t}.
\]
The odd part of $p_w(z^{-1})^2/(g(z^{-2})-wz)$ is
$wzS_t(z^2)$. Consequently the corner identity gives
$M_m(t)=[s^m]S_t(s)$ up to its finite inverse-row error.
\end{proposition}
\begin{proof}
Multiply numerator and denominator by $g(z^{-2})+wz$ and
retain the odd powers of $z$. The exact row-truncation identity
in Lemma~\ref{lem:corner-scattering-factor} supplies the error.
The next lemma bounds this error on the whole closed domain,
including $w=0$.
\end{proof}

\begin{lemma}[The scattering approximation on the whole closed domain]
\label{lem:closed-domain-scattering}
For each fixed polynomial background in the preceding small
neighborhood, there are $C>0$ and $0<\kappa<1$ such that
\[
 \sup_{t\in\overline D}|M_m(t)-[s^m]S_t(s)|\le C\kappa^m,
\]
where $S_t$ is the generating function in
Proposition~\ref{prop:fixed-background-forcing}. The same bound
holds for every Taylor coefficient after composition with $f$.
\end{lemma}
\begin{proof}
Use the scaled finite-section radius $\rho>1$ from
Proposition~\ref{prop:boundary-polynomial-stability}.
Choose $1<r_0<R^2<\rho$ inside its squared-variable collar.
For $t$ in the closed interior of $h(r_0\T)$, the far-factor
function $F_w$ is holomorphic on $|z|<R^{-1}$ and uniformly bounded
on that circle: any uncancelled physical zero has $|z|\le\sqrt{r_0}$
in the original exterior variable. Thus $|f_j|\le CR^j$.
The row truncation argument now gives low-row errors bounded by
$C(R/\rho)^N$. In the exact corner identity in
Lemma~\ref{lem:corner-scattering-factor}, multiplication by the
finitely many tail coefficients bounds the scattering error by
$C(R^2/\rho)^N$.

This error is holomorphic in $w$ on a fixed neighborhood of
$\mathcal W$. For even $N$ it is odd in $w$. Division by $w$
therefore preserves an exponential bound on $\mathcal W$: near zero
use Cauchy's estimate on a fixed small circle, and elsewhere divide
by a fixed lower bound for $|w|$. With $N=2m+2$, the parity formula
for $S_t$ gives the assertion, including $t=-1$. The coefficient
claim follows from the supremum bound on the unit disk after composition.
\end{proof}


\begin{lemma}[Coefficients of a Cauchy pole approaching the contact point]
\label{lem:moving-cauchy-coefficients}
Suppose $f:\mathbb D\to D$ extends to a $C^{1,\beta}$
diffeomorphism of the closed disk for some $\beta>0$, is bi-Lipschitz,
and satisfies $f(-1)=-1$, $f'(-1)>0$.
If $a\in D$ has $|a+1|\le C_0x$ and
$\operatorname{dist}(a,\partial D)\ge c_0x$, with $0<x<x_0$,
then
\[
 \left|[u^k]\frac1{a-f(u)}\right|\le C e^{C_1kx}.
\]
The same estimate, multiplied by $C/x$, bounds the derivative of
this coefficient with respect to $a$ on a smaller such region.
It also holds with a numerator $H(f(u))$, uniformly for a bounded
family of functions $H$ holomorphic on a fixed neighborhood of
$\overline D$. The exponential constant $C_1$ depends only on the
bi-Lipschitz bounds and $C_0,c_0$, not on the H\"older seminorm.
\end{lemma}
\begin{proof}
Put $v=f^{-1}(a)$. The bi-Lipschitz bounds imply
$|v+1|\le Cx$ and $1-|v|\ge cx$.
Remove the simple pole:
\[
 \frac1{a-f(u)}=\frac1{f'(v)(v-u)}+H_a(u).
\]
The remainder is holomorphic in the disk and is bounded there by
$C|u-v|^{\beta-1}$, by the H\"older Taylor remainder and the lower
Lipschitz bound. Its boundary $L^1$ norm is uniformly bounded,
so every Taylor coefficient is bounded by $C$. The pole term has
coefficient $1/(f'(v)v^{k+1})$, bounded by $Ce^{C_1kx}$.
Cauchy's estimate in the parameter $a$, on a disk of radius $cx$
contained in $D$, gives the derivative bound, with a possibly larger
$C_1$. To include a numerator, divide
$H(t)=H(a)+(t-a)\widetilde H_a(t)$; the divided differences are
uniformly bounded and holomorphic near $\overline D$.
\end{proof}


\begin{proposition}[The growing conformal forcing bound]
\label{prop:conformal-forcing-bound}
For every fixed polynomial background in the preceding small
neighborhood there are $\theta_0>0$ and $C_P$ such that
\begin{equation}\label{eq:conformal-forcing-coefficients}
 |[u^k]M_m(f(u))|\le C_Pm^{-3/2},
 \qquad 0\le k\le\theta_0m.
\end{equation}
Moreover $\theta_0$ can be chosen uniformly on a sufficiently small
fixed $\mathcal W^{1+\alpha}$ neighborhood of the squared unit
circle, for any fixed $\alpha>0$. The constants $C_P$ and the
threshold for $m$ may still depend arbitrarily on the fixed background.
Consequently, for $-1/2<s<0$ and fixed $0<\theta<\theta_0$,
\[
 \|P_{\lfloor\theta m\rfloor}(M_m\circ f)\|_{H^s}
       \le C_{P,s}m^{s-1}.
\]
\end{proposition}
\begin{proof}
Lemma~\ref{lem:quantitative-conformal-coordinate} gives a
$C^{1,\alpha}$ normalized map with $\|f'-1\|_\infty<1/4$ and
uniform bi-Lipschitz bounds. Real symmetry gives $f(-1)=-1$ and
$f'(-1)>0$.

It is enough, by Lemma~\ref{lem:closed-domain-scattering}, to bound
the double coefficient of $S_{f(u)}(s)$. Continue it from the $s$ disk
to a thin slit annulus $1\le|s|\le1+\delta$, cut along
$[-1-\delta,-1]$. Its only nonanalyticity there is the square root.
The continued function $h(s^{-1})$ remains within $C\delta$ of
$\Gamma$ on this annulus: its derivative is continuous across the
circle away from the cut and has a continuous limiting value at the
contact point on both sides. By decreasing $\delta$ for the fixed
background, the derivative bound here can be chosen absolute.
Choose $R=1-L\delta$, with a fixed sufficiently large $L$ depending
only on the bi-Lipschitz constants of $f$. Then $f(\{|u|\le R\})$
stays farther than $C\delta$ from $\Gamma$. No denominator zero
is crossed in the slit annulus for these $u$.

Deforming the $s$ coefficient contour gives an outer-circle term
and an integral of the two boundary values on the cut. The outer
term has $u^k$ coefficient at most
\[
 C_P(1+\delta)^{-m}R^{-k}.
\]
This is exponentially small if $k\le\theta_0m$, with
$\theta_0<1/(4L)$ and $\delta$ sufficiently small.

On the cut write $s=-1-x$, $\zeta=-1/(1+x)$, and
\[
 g_\pm=P(\zeta)\pm i\sqrt x,\qquad
 t_\pm=\zeta g_\pm^2-1.
\]
Since $P(\zeta)=O_P(x)$, these points satisfy
\[
 t_\pm=-1+x+O_P(x^{3/2}),\qquad
 |t_+-t_-|\le C_Px^{3/2}.
\]
They and the segment between them lie in $D$ at distance comparable
to $x$ from its boundary, after decreasing $\delta$ once more.
The two numerators in $S_t$ are
$N_\pm(t)=\zeta A_t(\zeta)+\zeta^2B_t(\zeta)g_\pm$.
They are uniformly bounded holomorphic functions of $t$ near
$\overline D$, and $N_+-N_-=O_P(\sqrt x)$ in that norm.
Split their jump as
\[
 \frac{N_+(t)-N_-(t)}{t_+-t}
 +N_-(t)\left(\frac1{t_+-t}-\frac1{t_--t}\right).
\]
Lemma~\ref{lem:moving-cauchy-coefficients}, and its parameter
derivative on the segment between $t_-$ and $t_+$, bound the
$u^k$ coefficient of this jump by
$C_P\sqrt x\,e^{C_1kx}$. This extraction is performed on a circle
$|u|=R$ smaller than both pole radii, as ensured by the choice of $L$.
The cut contribution is therefore bounded by
\[
 C_P\int_0^\delta x^{1/2}(1+x)^{-m-1}e^{C_1kx}\,dx
 \le C_Pm^{-3/2}
\]
when $k\le\theta_0m$ and $\theta_0<1/(4C_1)$.
The small circle about the branch point contributes zero in the
contour deformation, since the boundary values are bounded there
for $|u|\le R$. This proves \eqref{eq:conformal-forcing-coefficients}.

Only the uniform bi-Lipschitz constants and the fixed small
$C^1$ neighborhood enter the restrictions on $\theta_0$.
Lemma~\ref{lem:quantitative-conformal-coordinate} supplies these
uniformly on the stated $\mathcal W^{1+\alpha}$ neighborhood.
The slit-annulus width may
shrink with the particular polynomial; the restrictions on
$\theta_0$ above do not depend on that width. Finally, summing
$(k+1)^{2s}m^{-3}$ over $k<\theta m$ proves the Sobolev estimate.
\end{proof}


\begin{corollary}[Explicit parameters]
\label{cor:explicit-neighborhood-parameters}
In Proposition~\ref{prop:two-level-isomorphism} one may use
$\gamma=2^{-1000}$ and bound both operator norms and both inverse
norms by $2^{100}$. The forcing proportion may be taken as
$\theta_0=2^{-20}$. Thus
\[
 s=-\tfrac14,\qquad \tau=\tfrac18,\qquad \sigma=\tfrac38,
 \qquad \theta=2^{-10000}
\]
satisfy the strict oversampling inequality, with left side less
than $1/2$. These choices are intentionally inefficient.
\end{corollary}
\begin{proof}
We record a loose constant ledger to make the parameter choices
checkable. Use the coefficient norms already specified, and the
periodic H\"older norm
$\|v\|_\infty+\|v'\|_\infty+[v']_\alpha$ for $C^{1,\alpha}$.
All the following estimates apply at $r=3/4,7/8$ and $\alpha=1/8$.

In the endpoint lemma, $A_n\le(n+1)^{-1/2}$ gives Schur row
sums at most $24$ and column sums at most $28$, by splitting
the displayed sums at $j=k+1$. Thus $\|E\|\le64$ is valid.
The triangular inverse has squared Hilbert--Schmidt norm at most
$\sum_{\ell\ge1}\ell^{-3/2}(1+\log\ell)<8$, so its norm is
at most $8$. In the Cauchy-projection lemma the $L^2$ difference
is at most $4\delta$; differentiating gives at most $14\delta$
in the equivalent derivative norm. Interpolation and the Fourier
weight comparison give $64\delta$ on $H^r$. The off-diagonal
smoothing constant can be taken as $16$ by its two Neumann
resolvent factors.

The product estimates in the proof above give, successively,
\[
 \|(h-s)_+\|_{\mathcal W^{r+1}}\le20\gamma,\quad
 \|h'-1\|_{\mathcal W_{\rm pos}^r}\le21\gamma,\quad
 \|R-1\|_{\mathcal W_{\rm pos}^r}\le5\gamma,\quad
 \|m-1\|_{\mathcal W_{\rm pos}^r}\le52\gamma.
\]
For example, the three positive terms in $h-s$ cost at most
$16\gamma+12\gamma^2$. The direct, positive-density, and
negative-density errors in $B-I$ are bounded by
$3328\gamma$, $16384\gamma$, and $40960\gamma$ respectively.
The still looser bound $2^{25}\gamma$ therefore suffices.

The remaining norm comparisons can also be bounded generously.
For the two values of $r$, the elementary sine inequalities give
\[
 n^{2r}\le
 \int_{-\pi}^{\pi}
  \frac{|e^{int}-1|^2}{|e^{it}-1|^{1+2r}}\,dt
 \le64n^{2r}\qquad(n\ge1).
\]
Splitting the upper integral at $1/n$ proves the upper bound;
integrating only up to $1/n$ proves the lower bound.
Thus fractional composition by a circle map with both metric
and Jacobian bounds at most $2$ costs less than $2^{20}$.
For $a=-2s\in\{1/2,1/4\}$ the monomial beta weight lies between
$\tfrac12(k+1)^{-a}$ and $8(k+1)^{-a}$: for the lower bound
integrate over $0\le1-t\le1/(k+1)$, and for the upper use
$(1-x)^k\le e^{-kx}$. Analytic multiplier norms consequently
cost at most $8$ times the supremum norm.

For the conformal coordinate, radial conversion gives
$\|\ell\|_{C^{1,\alpha}}\le2^{12}\gamma$.
This follows by first taking the logarithm of
$h(e^{it})e^{-it}$ and then inverting its angle, whose derivative
is between $1/2$ and $3/2$. The cotangent-kernel proof bounds
the Hilbert transform in the two required H\"older norms by
$2^{20}$. In the fixed-point proof one may take
$M=2^{24}\|\ell\|_{C^{1,\alpha}}$; composition costs at most
$6\|\ell\|_{C^{1,\alpha}}$ on this set. It follows that
$\|f'-1\|_\infty\le2^{40}\gamma$ for $\gamma\le2^{-50}$.
In particular all the metric and Jacobian bounds used above
are at most $2$ for the chosen $\gamma$.

The semigroup proof gives
$\|b_xb_y/(b_x+b_y)\|_{\mathcal W}\le8$ and the associated
base two-variable kernel norm at most $20$.
The triangular-indicator argument has bound
$4n(1+\log n)\le16n^{1+\alpha}$.
Thus $2^{16}$ bounds both divided-difference constants
needed here, including positive weights up to order $2$.
Neumann series, the logarithm and exponential, the derivative
coefficient count, and the path of length at most $2|w|$
then bound the weighted sum for $A_t-1,B_t$ by $2^{30}\gamma$.

Combining the preceding bounds in the explicit factorization,
both bare norms are below $2^{40}$, and the general differentiated
density bound is below $2^{45}$ times its multiplier norm.
The actual perturbation, including all coefficient multipliers,
is below $2^{100}\gamma$. Its product with the bare inverse is
therefore less than $1/2$ at $\gamma=2^{-1000}$.
The two actual norms and inverse norms are consequently below
$2^{100}$.

Finally the forcing proof can use $R=1-32\delta$ for its
coefficient-extraction circle. After shrinking the collar for
the fixed background its continued derivatives are bounded by
$2$. The moving cut poles satisfy $|f^{-1}(a)+1|\le4x$;
also after the parameter Cauchy estimate their coefficient growth
can be bounded by $C_Pe^{64kx}$. These choices permit
$\theta_0=2^{-20}$. The collar width and prefactor may depend
on the background, as before. Now
\[
 2^{200}(4\cdot2^{-10000})^{1/8}<\tfrac12
\]
proves the stated oversampling assertion.
\end{proof}


\begin{lemma}[Oversampling from two Sobolev bounds]
\label{lem:sobolev-oversampling}
Suppose $A$ is an isomorphism on $H^s$ and $H^{s+\tau}$, with
$\tau>0$ and compatible inverses. Write
$M_s=\|A\|_{H^s\to H^s}$ and
$K_r=\|A^{-1}\|_{H^r\to H^r}$. If $q\ge h$ and
\[
 \varepsilon=M_sK_{s+\tau}
       \left(\frac h{q+1}\right)^\tau<1,
\]
then $P_hAP_q$ has a right inverse with $H^s$ norm at most
$K_s/(1-\varepsilon)$, whose output has degree less than $q$.
No small bound for $A-I$ is assumed.
\end{lemma}
\begin{proof}
On $P_hH^s$, set $E=P_hA(I-P_q)A^{-1}P_h$.
The tail estimate
$\|(I-P_q)v\|_{H^s}\le(q+1)^{-\tau}\|v\|_{H^{s+\tau}}$
and the finite-band estimate
$\|P_hv\|_{H^{s+\tau}}\le h^\tau\|P_hv\|_{H^s}$
give $\|E\|\le\varepsilon$. Since
$P_hAP_qA^{-1}P_h=I-E$, the operator
$P_qA^{-1}P_h(I-E)^{-1}$ is the required right inverse.
\end{proof}


\begin{proposition}[Two explicit analytic estimates would suffice locally]
\label{prop:sobolev-completion-reduction}
Fix one polynomial background $P$ as above and numbers
$0<\theta<\sigma<1/2$, $-1/2<s<0$, and $\tau>0$.
Let $f$ be the exact conformal map onto $D$, with $f(0)=0$ and
$f'(0)>0$, and define on polynomials
\[
 \mathcal A_gx=(\mathcal J_g((J^0)^{-1}x))\circ f.
\]
Assume the following two estimates, established above on the
specified small neighborhood and with the explicit parameters.
\begin{enumerate}
\item $\mathcal A_g$ extends to compatible isomorphisms on $H^s$
and $H^{s+\tau}$, with
\[
 \|\mathcal A_g\|_{H^s}\,
 \|\mathcal A_g^{-1}\|_{H^{s+\tau}}
            (\theta/\sigma)^\tau<1.
\]
\item For $h=\lfloor\theta m\rfloor$, the exact forcing satisfies
\[
 \|P_h(M_m\circ f)\|_{H^s}\le C_Pm^{s-1}.
\]
\end{enumerate}
Then, for all sufficiently large $m$, real corrections supported at
$g_{-j}$, $m-\lfloor\sigma m\rfloor+1\le j\le m$, make the first
$h$ jets vanish exactly. Their preconditioned vectors satisfy
$\|x\|_{H^s}=O_P(m^{s-1})$, so all the coefficient and geometry
conclusions of Corollary~\ref{cor:negative-sobolev-correction} hold.
\end{proposition}
\begin{proof}
Lemma~\ref{lem:sobolev-oversampling} gives uniformly bounded right
inverses for the first $h$ rows and $q=\lfloor\sigma m\rfloor$
columns of the preconditioned limiting operator. By
Corollary~\ref{cor:boundary-jacobian-replacement}, the actual finite
linear maps differ in $H^s$ operator norm by an exponentially small
quantity times a fixed power of $m$. Indeed
\eqref{eq:base-preconditioner} bounds the input coefficient norm by
$C_sq^{-s}\|x\|_{H^s}$, while $h$ uniformly bounded output
coefficients have $H^s$ norm at most $C_sh^{s+1/2}$ times their
maximum. A Neumann correction therefore supplies uniformly bounded
right inverses $B_m$ for the finite maps.

Write the exact finite equations as $b_m+A_mx+N_m(x)=0$.
Corollary~\ref{cor:conformal-feedback-bound} and the same norm
conversions bound the nonlinear remainder, and its difference on
a ball, by an exponentially small quantity times fixed powers of
$m$ and the corresponding quadratic expression in $x$. The
difference estimate also follows directly by subtracting the
convergent separated-corner resolvent series. On a ball of radius
$C m^{s-1}$, the map
\[
 x\longmapsto-B_mb_m-B_mN_m(x)
\]
therefore maps the ball into itself and has Lipschitz constant less
than $1/2$ for large $m$, after increasing the fixed $C$.
Its fixed point solves the exact equations. Since the background
and the normalized conformal map have real coefficients, all these
maps preserve real coefficient sequences, and the fixed point is real.
The odd-pencil identity turns vanishing jets into the asserted
characteristic-polynomial divisibility. Invisible restoration and
the preceding corollary give the final bounds.
\end{proof}


\section{Exact selection and the infinite symbol}
The following gives an exact selection rule. Its purpose is to remove
any dependence on an unspecified choice of a nonlinear root. The
algorithm is not intended for practical computation.

\begin{theorem}[A continuous counterexample]
\label{thm:counterexample}
There is a computably specified continuous complex scalar symbol $a$
with neither separate one-sided annular holomorphic extension and an
increasing sequence $n_j$ such that $T_{n_j}(a)$ has both $1$ and $-1$
as eigenvalues of algebraic multiplicity at least
$\lfloor\theta(n_j-1)/2\rfloor$, where $\theta=2^{-10000}$.
Both points are outside $a(\T)$. In particular its empirical
eigenvalue measures do not converge to $a_*m$.
\end{theorem}
\begin{proof}
Fix
\[
 \gamma=2^{-1000},\quad \theta=2^{-10000},\quad
 \sigma=\tfrac38,\quad m_0=2^{10000},\quad
 \kappa=2^{-1010},\quad \beta_j=\gamma\,2^{-j-4}.
\]
Start with $P_-=P_+=0$. At stage $j$, with the last chosen order
$m_{j-1}$, first add the positive packet
\[
 Q_j(s)=\tau_j s^{m_{j-1}}(1+s),\qquad
 \tau_j=\kappa\,2^{-\lceil\sqrt{2m_{j-1}+1}\rceil}.
\]
We will always choose $m_j\ge8m_{j-1}$, divisible by eight.
The new packet vanishes at $-1$.

All positive costs fit their stage budgets. Its $\mathcal W^4$ norm
is bounded by $C\tau_j(1+m_{j-1})^4$. The same bound, with a fixed
larger constant, bounds its change to the squared curve in
$\mathcal W^{1+\alpha}$, $\alpha=1/8$. Indeed
$g_0(1+s)$ belongs to that algebra, and the old finite polynomial
has support of magnitude at most $2m_{j-1}+1$ and bounded
$\mathcal W^0$ norm. Expand the square and use the weighted
convolution inequality. The constant $2^{100}$ is more than
sufficient in these estimates. For $x\ge m_0$ the function
$(1+x)^4 2^{-\sqrt{2x+1}}$ decreases and its value at $8x$
is less than a quarter of its value at $x$. Its initial value
is smaller than the required bound for $\beta_1$ by the displayed
choices. Thus each positive $\mathcal W^4$ cost and squared-curve
cost is strictly less than $\beta_j$, uniformly in the subsequent
choices of orders.

Hold the resulting background $P$ fixed. For each prospective
$m\ge8m_{j-1}$ divisible by eight put
$q=3m/8$ and $h=\lfloor\theta m\rfloor$.
Let $v=(v_0,\ldots,v_{q-1})$ be real and set
\[
 L_v(s)=s^{-m}\sum_{d=0}^{q-1}v_ds^d,\qquad
 D_v(s)=L_v(s)-L_v(-1)K_m(s).
\]
Use the exact $K_m$ in Lemma~\ref{lem:localized-restoration}.
Require the first $h$ coefficients in $t=w^2-1$ of
\[
 \frac{\det(wI-T_{2m+1}(z(g+L_v)(z^2)))}{w}
\]
to vanish. These are $h$ real polynomial equations in the $q$
variables $v_d$. There is no division ambiguity: the odd Toeplitz
symbol makes the characteristic polynomial odd, and its quotient
by $w$ is a polynomial in $w^2$, hence in $t$.
The coefficients of these equations are computable from the
current finite background and the rational binomial coefficients.

Put $x=J^0v$ and impose also the strict inequalities
\[
 R_1=\sum_{k<q}\frac{|x_k|}{k+1}<m^{-19/16},\qquad
 R_{1/2}=\sum_{k<q}\frac{|x_k|}{\sqrt{k+1}}<m^{-15/16}.
\]
Require, separately, each of the following increments to be
strictly below $\beta_j$:
\[
 \|D_v\|_{\mathcal W^0},\qquad
 \|D_v/g_0\|_{\mathcal W^0},\qquad
 \|g_0D_v\|_{\mathcal W^{1+\alpha}},\qquad
 \|s(g+D_v)^2-sg^2\|_{\mathcal W^{1+\alpha}}.
\]
The last norm may be certified by the triangle bound obtained
from $2sg_0D_v+2sPD_v+sD_v^2$.

These conditions have regular solutions for all sufficiently
large $m$. The two-level theorem, forcing bound, and oversampling
lemma give the local completion with
$\|x\|_{H^{-1/4}}=O_P(m^{-5/4})$.
Thus $R_1=O_P(m^{-5/4})$ and $R_{1/2}=O_P(m^{-1})$, which
satisfy both strict bounds eventually. Conversely the stated
strict bounds alone imply, by the restoration estimates,
\[
 \|D_v\|_{\mathcal W^0}=O(m^{-15/16}),\quad
 \|g_0D_v\|_{\mathcal W^{1+\alpha}}=O(m^{-1/16}),\quad
 \|D_v/g_0\|_{\mathcal W^0}
       =O(m^{-3/16}\sqrt{\log(m+1)}).
\]
For the third bound use $H^{3/4}(\T)\subset\mathcal W^0$.
The squared-curve increment is bounded by
$C_P(m^{-1/16}+m^{-3/4})$. Hence the open norm budgets hold
eventually as well.

The polynomial equations have full row rank at the supplied
solutions. In the local completion proof their Jacobian differs
by an exponentially small operator from a map with uniformly
bounded right inverse. At a zero jet, multiplication by
$\det(H-tU)$ and composition with $f$ induce invertible
triangular maps on the first $h$ jets. The original polynomial
equations therefore have the same full row rank.

Here is the deterministic selection rule. Enumerate all finite
rational encodings consisting of an admissible
$m$, a subset of $h$ coordinates, rational values for the other
coordinates, a rational cube $B(c,r)$ in the selected coordinates,
and a rational square matrix $K$. If $F$ is the resulting square
polynomial system, seek strict interval certificates
\[
 \sup_{B(c,r)}\|I-KDF\|_\infty<\tfrac12,\qquad
 \|KF(c)\|_\infty<r/2.
\]
Also certify both $R$ bounds and all four norm budgets throughout
the cube. To fix the order completely, write each rational in
reduced form with positive denominator, list the coordinate subset
in increasing order, and list matrix entries in row order. The height
of an encoding is the maximum of $m$ and the absolute values of all
its integer numerators and denominators. In round $N$ inspect all
encodings of height at most $N$ in lexicographic order, with absolute
coefficient-enclosure error at most $2^{-N}$, and select the first
one whose strict certificates succeed. Each round is finite.
Use tail-enclosure tolerance at most $\min(2^{-N},\beta_j/8)$.
The displayed estimates make $z\mapsto z-KF(z)$ a contraction
of the cube into itself. They also imply that $K$ is invertible.
Its unique fixed point is the selected exact root, and the
contraction iteration specifies that root as a computable real vector.

The coefficients remain computable by induction. The explicit stage
budgets also give a computable uniform modulus of convergence for
the final Fourier series. All norm certificates are computable. For a restored polynomial
of maximum degree magnitude $N_0$, the exact equality $D_v(-1)=0$
and the binomial recurrences give, for $n>2N_0$,
\[
 |[s^{-n}]g_0D_v|\le 8N_0\|D_v\|_{\mathcal W^0}n^{-5/2},
 \qquad
 |[s^{-n}]D_v/g_0|\le 8N_0\|D_v\|_{\mathcal W^0}n^{-3/2}.
\]
Finite sums and integral tail bounds therefore give upper
enclosures of the two infinite norms. More explicitly, past
$J>2N_0$ their respective omitted norm tails are at most
$128N_0\|D_v\|_{\mathcal W^0}J^{-3/8}$ and
$16N_0\|D_v\|_{\mathcal W^0}J^{-1/2}$.
The constants follow from $A_n\le(n+1)^{-1/2}$ and the
first-difference recurrences for $A_n$ and $\binom{1/2}{n}$.
Use an enclosure tolerance
less than a quarter of the stage budget. The remaining terms
are finite polynomial coefficient norms. The cancellation in
$D_v$ is structural and holds throughout every cube, so this
tail estimate does not use an approximately zero endpoint value.

The search terminates. At a regular root supplied above some
$h$ columns of the Jacobian form an invertible matrix. By the
implicit-function theorem the other coordinates may be chosen
rationally nearby while preserving a root and the strict budgets.
A sufficiently small rational cube and a sufficiently accurate
rational inverse preconditioner satisfy both contraction
inequalities. Refining the coefficient and tail enclosures
eventually certifies all the strict conditions. The dovetailed
search reaches such a certificate in a finite round.

Choose this $m$ as $m_j$, add its exact $D_v$ to $P_-$, and
continue. The sums of the positive and negative stage costs are
smaller than $\gamma/8$. Thus all finite backgrounds remain
strictly within the neighborhood of the operator theorem:
the three negative norms and the positive $\mathcal W^4$ norm
are controlled separately, and both types of squared-curve
increments have summable bounds. This proves the induction.

The resulting series converges in $\mathcal W^0$ to a continuous
$g$, hence to a continuous $a(z)=zg(z^2)$.
Every later change is invisible to a previously selected
$T_{2m_j+1}(a)$. The next positive packet starts at Fourier
frequency $2m_j+1$ in $a$; the next negative correction starts
beyond $4m_j$ in the $g$ variable; and the current restoring
packet starts at $g_{-m_j-1}$. All lie outside the diagonals
used by the previously fixed section. Its exact characteristic
factor $(w^2-1)^{\lfloor\theta m_j\rfloor}$ therefore persists
in the final symbol.

The positive coefficient at $g_{m_{j-1}}$ is exactly $\tau_j$,
and $\log\tau_j/(2m_{j-1}+1)\to0$. At $g_{-3m_j}$ no
correction ever acts: the current restoration ends at
$-2m_j-1$ and the next correction starts beyond $-4m_j$.
These coefficients are the original nonzero binomial coefficients,
of order $m_j^{-3/2}$. Consequently both positive and negative
Fourier root limsups of $a$ equal one. An outer annular extension
would give $|a_k|\le C_\rho\rho^{-k}$ for some $\rho>1$ by
Cauchy's Laurent coefficient formula; an inner extension would
give the analogous geometric decay of $a_{-k}$. Both are excluded.

Finally $a_0(\T)$ satisfies $|w^2-1|=1$ and its distance from
either $1$ or $-1$ is $\sqrt2-1$. For example,
$1=|w-1||w+1|\le d(2+d)$, where $d=|w-1|$,
and equality is attained at $w=\sqrt2$.
The total uniform perturbation is smaller than $2\gamma$, so
both points remain outside $a(\T)$. The explicit test function
$\Phi(w)=\max\{0,1-8|w-1|\}$ is continuous, nonnegative,
compactly supported, and zero on $a(\T)$.
Its integral against $a_*m$ is zero. At $n_j=2m_j+1$ its empirical eigenvalue
integral is at least
$\lfloor\theta m_j\rfloor/(2m_j+1)\to\theta/2>0$.
This proves the asserted failure of canonical convergence.
\end{proof}

\begin{thebibliography}{9}
\bibitem{BBG} J. M. Bogoya, A. B\"ottcher, S. M. Grudsky,
Asymptotics of individual eigenvalues of a class of large Hessenberg
Toeplitz matrices, OTAA \textbf{220} (2012), 77--95.
\url{https://doi.org/10.1007/978-3-0348-0346-5_5}.
Author manuscript: \url{https://www.math.cinvestav.mx/~grudsky/Papers/116.pdf}.

\bibitem{BFGM} A. B\"ottcher, L. Fukshansky, S. R. Garcia, H. Maharaj,
Toeplitz determinants with perturbations in the corners,
J. Funct. Anal. \textbf{268} (2015), 171--193.
\url{https://doi.org/10.1016/j.jfa.2014.10.023}.
Section 3, equation (6), pp. 174--175, states the finite-section
stability result used here. The continuous nonvanishing
zero-index case is within its hypotheses.

\bibitem{Clopen} clopen.solutions,
\emph{OP-00159: Widom's canonical distribution conjecture for Toeplitz
eigenvalues}. Online problem catalogue.
\url{https://clopen.solutions/OP-00159/}.
Accessed 7 October 2026.

\bibitem{OpenProblemsInNLA} OpenProblemsInNLA contributors,
\emph{SP-14: Widom's canonical distribution conjecture for Toeplitz
eigenvalues}. Open Problems in Numerical Linear Algebra,
GitHub repository (2026).
\url{https://github.com/ajt60gaibb/OpenProblemsInNLA/blob/main/eigenvalues-and-inverse-problems/SP-14/README.md}.
Accessed 9 October 2026.
\end{thebibliography}

\section*{AI usage disclosure}
OpenAI's Codex was used extensively to develop the mathematical
construction and arguments, search and check references, write and run
computational experiments, draft this manuscript, and conduct an
internal review of the argument. This AI-assisted review does not
constitute independent peer review or formal proof-assistant verification.
\end{document}
